\documentclass[11pt,a4paper]{article}
% BEGIN INCLUDED SOURCE: preamble
\usepackage[T1]{fontenc}
\usepackage{lmodern,microtype}
\usepackage[a4paper,margin=25mm,headheight=15pt]{geometry}
\usepackage{amsmath,amssymb,amsthm,mathtools}
\usepackage{booktabs,longtable,array,enumitem}
\usepackage[dvipsnames]{xcolor}
\usepackage{fancyhdr}
\usepackage[colorlinks=true,linkcolor=MidnightBlue,citecolor=MidnightBlue,
 urlcolor=MidnightBlue,breaklinks=true]{hyperref}
\usepackage{bookmark}
\pagestyle{fancy}
\fancyhf{}
\lhead{\small Gamma determinants and colored moments}
\rhead{\small ProveIt research continuation}
\cfoot{\thepage}
\renewcommand{\headrulewidth}{0.3pt}
\numberwithin{equation}{section}
\newtheorem{theorem}{Theorem}[section]
\newtheorem{lemma}[theorem]{Lemma}
\newtheorem{proposition}[theorem]{Proposition}
\newtheorem{corollary}[theorem]{Corollary}
\theoremstyle{definition}
\newtheorem{definition}[theorem]{Definition}
\theoremstyle{remark}
\newtheorem{remark}[theorem]{Remark}
\DeclareMathOperator{\Li}{Li}
\DeclareMathOperator{\FP}{FP}
\DeclareMathOperator{\CT}{CT}
\DeclareMathOperator{\Var}{Var}
\DeclareMathOperator{\sech}{sech}
\DeclareMathOperator{\NF}{NF}
\newcommand{\dd}{\,\mathrm d}
\newcommand{\ii}{\mathrm i}
\newcommand{\N}{\mathbb N}
\newcommand{\R}{\mathbb R}
\newcommand{\C}{\mathbb C}
\newcommand{\E}{\mathbb E}
\newcommand{\Q}{\mathbb Q}
\newcommand{\src}[1]{\nolinkurl{#1}}
\newcommand{\pin}{\src{c4036d60715916ae72dd18983782fdb182840b86}}
\setlist{itemsep=4pt,topsep=6pt,leftmargin=*}
\setlength{\emergencystretch}{2em}
\allowdisplaybreaks[2]
\hypersetup{
 pdftitle={Gamma Determinants and Colored Stieltjes Moments},
 pdfauthor={Research continuation prepared for Vladimir Reshetnikov with OpenAI assistance},
 pdfsubject={Exact Gamma normalizations, finite Gaussian trigamma moments, colored Lerch moments, and clustered polylogarithms}
}

% END INCLUDED SOURCE: preamble

\title{\textbf{Gamma Determinants and\\Colored Stieltjes Moments}\\[7pt]
 \large Closed normalization constants, finite Gaussian moment laws,\\
 and clustered polylogarithm identities}
\author{Research continuation prepared for Vladimir Reshetnikov's ProveIt programme\\
 \small Mathematical derivations and independent computational checks with OpenAI assistance}
\date{12 October 2026}
\begin{document}
\maketitle
\begin{abstract}
We resolve the incoming Gamma-basis normalization question at every
dimension by reducing its coefficient determinant to Gustafson's
classical multivariable Barnes integral. The resulting constants are
explicit rational multiples of powers of $\pi$, with a two-step
recurrence. A parameter deformation yields a further exact phenomenon:
under an explicit positive $n$-variable Gamma measure, a normalized sum
of trigamma values has precisely the first $n+1$ moments of a standard
Gaussian. This holds at every positive Gamma shift and implies a
Gaussian limit even when that shift varies with the dimension.

In a separate translation geometry, we solve the colored analogue of
the incoming balanced Stieltjes-moment problem. A unique homogeneous
combination of Lerch order derivatives cancels its first inverse-power
tails. In the general convergence range, each positive-integer moment reduces to one
negative-order Lerch value times a finite ordinary-zeta combination.
Root-of-unity colors give explicit generalized Stieltjes identities;
exceptional alternating moments evaluate to Dirichlet beta values.
The color tending to one has an exact, normalization-dependent
constant term, which we compute.

We also give a finite multiple-polylogarithm dictionary for Gamma
spectral faces with two groups of repeated masses, and an explicit
depth-two reduction for a printed Gaussian-color example.
Classical input identities, new consequences relative to the inspected
reports, numerical diagnostics, and remaining conjectures are kept
distinct. In particular, the Gaussian $S_6$ and revised $S_8$ conjectures
and the unrestricted unequal-tail prefix problem are not settled.
\end{abstract}
\vfill
{\noindent\small\textbf{Repository anchor.} The manuscript and all eleven
archives present in the supplied incoming directory were obtained from
\pin. Relevant mathematical passages were reviewed; this is not a
claim of a complete audit of the entire repository.\par}
\clearpage
{\small\tableofcontents}
\clearpage
% BEGIN INCLUDED SOURCE: sections/00_scope
\section{Problems addressed and mathematical status}
\label{sec:scope}

The current ProveIt manuscript develops polylogarithm identities
together with harmonic-tail, Gamma, Hurwitz, and generalized Stieltjes
calculi \cite{repo}. Recent incoming reports contain both explicit open
questions and extensive proved developments. A continuation must therefore
identify its exact starting point before claiming an advance.

The principal target here is Question 2 in the research section of
\emph{Ordered Boundary Identities and Gamma Kernels} \cite{repoordered}.
That report constructs a Gamma integral basis for a fixed-gap harmonic
recurrence and proves that its initial determinant equals
$(-1)^{m(m-1)/2}K_m$, where $K_m$ is a positive integral. It asks for
closed evaluations of $K_2,K_3,\ldots$. We give one formula for every
$m$, rather than a table of separately evaluated integrals.

The second target is Question 4 of \emph{Balanced Stieltjes Moments}
\cite{repobalanced}: a colored or Lerch analogue of the homogeneous
balancing and its polynomial moment reductions. The nontrivial color
removes the order pole, so the correct generator has no separate
$1/z$ subtraction. This changes the degree of the negative-order
coordinate and makes the limit of color one singular. We prove the
colored theorem on its full natural slit domain and give its exact
connection with the uncolored normalization.

The repeated-mass theorem answers a two-cluster part of the spectral
face question in \cite{repolattice}. It reduces arbitrary multiplicities
to finitely many multiple polylogarithms. The number of lattice terms
remains infinite: the finite dictionary describes the types and depths
of polylogarithms in each term, and does not give a finite evaluation
of the complete lattice sum.

\subsection{A concise result ledger}

\begin{longtable}{@{}>{\raggedright\arraybackslash}p{0.25\textwidth}>{\raggedright\arraybackslash}p{0.45\textwidth}>{\raggedright\arraybackslash}p{0.24\textwidth}@{}}
\toprule
Result & Exact advance & Status\\
\midrule
\endhead
Gamma normalization &
Closed $K_m$ for all $m$; explicit two-step recurrence &
Incoming Question 2 resolved\\[3pt]
Trigamma moments &
After normalization, the first $n+1$ moments are exactly Gaussian
for every positive shift &
Proved consequence of a parameter identity\\[3pt]
Colored balancing &
Unique coefficients in a specified homogeneous ansatz; all integer
moments and exceptional resonances &
Colored counterpart resolved in that ansatz\\[3pt]
Color-one limit &
Explicit pole subtraction and a shift-dependent constant-term
conversion &
Proved normalization identity\\[3pt]
Two mass groups &
Finite multiple-polylogarithm formula at every pair of multiplicities &
Two-cluster sector resolved\\[3pt]
Depth-two parity &
Exact double-shuffle and distribution certificate for one
printed example &
Classical parity made constructive\\[3pt]
$S_6$, revised $S_8$ &
No exact certificate supplied here &
Remain conjectural\\[3pt]
Unequal-tail prefix &
Initial basis normalization is evaluated; later determinant
nonvanishing is a separate issue &
Unrestricted problem remains open\\
\bottomrule
\end{longtable}

\subsection{Conventions and novelty}

The Gamma function and its logarithmic derivative are denoted by
$\Gamma$ and $\psi$; $\psi^{(j)}$ always means argument differentiation.
An order derivative of Hurwitz zeta or Lerch's transcendent is written
$\partial_s^j$. The generalized Stieltjes normalization is
\begin{equation}\label{eq:stieltjes-convention}
 \zeta(1+z,a)=\frac1z+
  \sum_{j=0}^{\infty}\frac{(-1)^j\gamma_j(a)}{j!}z^j,
 \qquad \gamma_0(a)=-\psi(a).
\end{equation}
All real-variable improper integrals asserted as ordinary integrals
converge absolutely. A constant term is always taken in a specified
parameter; no unspecified regularization is used.

Gustafson's integral, the Gamma reflection formula, the Lerch integral
representation, Bernoulli and Euler special values, and
multiple-polylogarithm Landen transformations are classical
\cite{gustafson1994,derkachov2017,dlmfgamma,dlmflerch,dlmfhurwitz,okudaueno}.
Our normalization theorem is a new application to the specific
incoming basis; the colored and clustered statements are explicit
extensions of the inspected incoming formulas. The finite trigamma
moment law is derived here from the complete parameter identity.
A targeted literature search supports these comparisons but does not
establish worldwide historical priority. No arithmetic independence,
minimal period basis, or proof-assistant formalization is claimed.


% END INCLUDED SOURCE: sections/00_scope

% BEGIN INCLUDED SOURCE: sections/01_gamma
\section{An exact Gamma moment determinant}
\label{gn:sec:determinant}

The incoming ordered-boundary report asks for closed evaluations of the
positive constants $K_m$ that normalize its Gamma integral basis
\cite[Section~8 and Section~9, question~2]{repoordered}.  All of these
constants can be evaluated.  The decisive step is to remove the source's
last moment by its already proved orthogonality relation and recognize
the remaining determinant as Gustafson's symplectic Barnes integral.
We first state the parameterized determinant with its full normalization.
This is a reformulation of a classical integral, and it is credited as
such; its application to the particular source constants and the moment
identities developed below are the results of the present continuation.

For $n\ge1$ and $\mathbf a=(a_1,\ldots,a_{2n+2})$ satisfying
$\Re a_\ell>0$, put
\begin{equation}\label{gn:eq:weight}
 W_{\mathbf a}(x)=\prod_{\ell=1}^{2n+2}
       \Gamma(a_\ell+ix/2)\Gamma(a_\ell-ix/2),
 \qquad A=\sum_{\ell=1}^{2n+2}a_\ell.
\end{equation}
Define the matrix and its determinant by
\begin{equation}\label{gn:eq:matrix}
 \mathsf B_n(\mathbf a)_{r,k}
 =\int_0^\infty x^{2r-1}\sinh(k\pi x)W_{\mathbf a}(x)\dd x,
 \qquad
 \mathcal D_n(\mathbf a)=\det\mathsf B_n(\mathbf a),
 \quad 1\le r,k\le n.
\end{equation}
The sinh columns are in increasing order.  This convention is relevant
when the result is applied to a source matrix with its columns reversed.

\begin{theorem}[Parameterized Gamma moment determinant]
\label{gn:thm:determinant}
For $\Re a_\ell>0$,
\begin{equation}\label{gn:eq:determinant}
 \boxed{\displaystyle
 \mathcal D_n(\mathbf a)
 =2^{2n^2}\pi^{n(n+1)}
       \frac{\prod_{\ell<j}\Gamma(a_\ell+a_j)}{\Gamma(A)}.}
\end{equation}
Every entry is an absolutely convergent ordinary integral.  The entries,
their determinant, and every fixed mixed parameter derivative are
holomorphic on this domain and may be differentiated under the integral.
The determinant is nonzero throughout the domain and is positive when
all parameters are positive real numbers.
\end{theorem}

\begin{proof}
Stirling's formula uniformly on compact parameter sets gives
$W_{\mathbf a}(x)=O(x^C e^{-(n+1)\pi x})$ for some finite $C$.
The largest sinh factor grows as $e^{n\pi x}$.  A fixed parameter
derivative adds at most powers of $\log(2+x)$ and preserves an integrable
exponential majorant.  At zero all Gamma arguments remain in a compact
subset of the right half-plane.  These observations prove the asserted
convergence and permit the uses of Fubini and parameter differentiation.

Write
\[
 \Delta(y_1,\ldots,y_n)=\prod_{j<k}(y_k-y_j)
\]
and introduce the integral
\begin{equation}\label{gn:eq:I}
\begin{split}
 \mathcal I_n(\mathbf a)
 ={}&\int_{(0,\infty)^n}
    \Delta(x_1^2,\ldots,x_n^2)
    \Delta(\cosh\pi x_1,\ldots,\cosh\pi x_n)\\
 &\hspace{12mm}\times
       \prod_{j=1}^n x_j\sinh(\pi x_j)W_{\mathbf a}(x_j)
          \dd x_1\cdots\dd x_n.
\end{split}
\end{equation}
Expanding the two determinants in the continuous Cauchy--Binet identity
and using Fubini gives
\begin{equation}\label{gn:eq:andreief}
 \mathcal D_n(\mathbf a)
      =\frac{2^{n(n-1)/2}}{n!}\mathcal I_n(\mathbf a).
\end{equation}
Indeed, the polynomial alternant is
$\det(x_j^{2r-1})=\prod_jx_j\,\Delta(x_1^2,\ldots,x_n^2)$,
while
\[
 \det(\sinh(k\pi x_j))_{j,k=1}^n
  =2^{n(n-1)/2}\prod_j\sinh(\pi x_j)
                   \Delta(\cosh\pi x_1,\ldots,\cosh\pi x_n).
\]
The second identity follows from
$\sinh(ku)=\sinh(u)U_{k-1}(\cosh u)$ and the leading coefficient
$2^{k-1}$ of the Chebyshev polynomial $U_{k-1}$.

For precision, the classical integral used next is Gustafson's
symplectic Barnes integral \cite{gustafson1994}, in the normalization
reproduced in \cite[equation~(2)]{derkachov2017}:
\begin{equation}\label{gn:eq:gustafson}
\begin{split}
 &\int_{\mathbb R^n}
 \frac{\displaystyle\prod_{j=1}^n\prod_{\ell=1}^{2n+2}
                 \Gamma(a_\ell+iu_j)\Gamma(a_\ell-iu_j)}
 {\displaystyle\prod_j\Gamma(2iu_j)\Gamma(-2iu_j)
       \prod_{j<k}\Gamma(iu_j+iu_k)\Gamma(iu_j-iu_k)
                    \Gamma(-iu_j+iu_k)\Gamma(-iu_j-iu_k)}
                  \dd\mathbf u\\
 &\hspace{22mm}=(2\pi)^n2^nn!
       \frac{\prod_{\ell<j}\Gamma(a_\ell+a_j)}{\Gamma(A)}.
\end{split}
\end{equation}
The right-half-plane hypotheses place all numerator poles on the
appropriate sides of the imaginary integration axes in its Barnes
form.  The integrand is even in each $u_j$, so the positive-orthant
integral equals $(2\pi)^nn!$ times the Gamma quotient.

Set $x_j=2u_j$.  Reflection yields the single-variable factor
\begin{equation}\label{gn:eq:single-reflection}
 \frac1{\Gamma(ix_j)\Gamma(-ix_j)}
          =\frac{x_j\sinh(\pi x_j)}\pi,
\end{equation}
and each four-factor interaction contributes
\begin{equation}\label{gn:eq:pair-reflection}
\begin{split}
 &\left\{
 \Gamma\left(\tfrac i2(x_j+x_k)\right)
 \Gamma\left(-\tfrac i2(x_j+x_k)\right)
 \Gamma\left(\tfrac i2(x_k-x_j)\right)
 \Gamma\left(-\tfrac i2(x_k-x_j)\right)
 \right\}^{-1}\\
 &\hspace{14mm}=
 \frac{(x_k^2-x_j^2)(\cosh\pi x_k-\cosh\pi x_j)}{8\pi^2}.
\end{split}
\end{equation}
The factor $2^{-n}$ from $\dd\mathbf u$ must also be retained.
Consequently \eqref{gn:eq:gustafson} gives
\begin{equation}\label{gn:eq:I-evaluation}
 \mathcal I_n(\mathbf a)
 =2^{2n+3n(n-1)/2}\pi^{n(n+1)}n!
       \frac{\prod_{\ell<j}\Gamma(a_\ell+a_j)}{\Gamma(A)}.
\end{equation}
Combining this with \eqref{gn:eq:andreief} proves
\eqref{gn:eq:determinant}, since the two powers of $2$ add to $2n^2$.
The Gamma quotient is nonzero in the stated domain.  For positive
parameters the two Vandermonde factors in \eqref{gn:eq:I} have the same
sign, and the integrand is strictly positive away from collisions.
This proves the positivity assertion as well.
\end{proof}

The case $n=0$ may be included with the empty determinant and empty
integral both equal to one.  Its Gamma quotient is then
$\Gamma(a_1+a_2)/\Gamma(a_1+a_2)=1$.

\section{Closed evaluation of every Gamma-basis normalization}
\label{gn:sec:normalization}

We recall exactly the source data needed for the application.  For an
integer $m\ge1$, set
\begin{align}
 W_m(x)&=|\Gamma(\tfrac14+ix/2)|^{4m},
       &\phi_r(x)&=\Im(x+i/2)^{2r},\label{gn:eq:source-weight}\\
 P_j(x)&=\sinh(\pi(m-j)x)\quad(1\le j<m),
       &P_m(x)&=\tanh(\pi x),\label{gn:eq:source-multipliers}\\
 t_{r,j}&=\frac4{(2r)!}\int_0^\infty
                 \phi_r(x)P_j(x)W_m(x)\dd x,
       &T_m&=(t_{r,j})_{r,j=1}^m.\label{gn:eq:source-matrix}
\end{align}
The incoming report proves the exact top moments
\begin{equation}\label{gn:eq:source-top}
 t_{m,j}=0\quad(j<m),\qquad
 t_{m,m}=(-1)^{m+1}\frac{2^{2m+2}\pi^m}{(2m)!},
\end{equation}
by Gamma recurrence and a contour shift crossing precisely two tanh
poles.  It also proves
\begin{equation}\label{gn:eq:source-normalization}
 \det T_m=(-1)^{m(m-1)/2}K_m,\qquad K_m>0.
\end{equation}
Equations \eqref{gn:eq:source-top} and
\eqref{gn:eq:source-normalization} are inherited results, specifically
the source theorem \texttt{gb:basis}, corollary
\texttt{gb:moment-identities}, and proposition
\texttt{gb:normalization} \cite{repoordered}.  The following evaluation
answers the open question about the remaining positive constant.

\begin{theorem}[All Gamma normalization constants]
\label{gn:thm:K}
For every integer $m\ge1$,
\begin{equation}\label{gn:eq:K}
 \boxed{\displaystyle
 K_m=\frac{2^{2m^2+2}(m-1)!\,\pi^{m(4m-1)/2}}
              {\Gamma(m/2)\prod_{r=1}^m(2r)!}.}
\end{equation}
In particular,
\begin{equation}\label{gn:eq:K-values}
 K_1=8\pi,\qquad
 K_2=\frac{64}{3}\pi^7,\qquad
 K_3=\frac{16384}{135}\pi^{16}.
\end{equation}
The constants satisfy the two-step recurrence
\begin{equation}\label{gn:eq:K-recurrence}
 \boxed{\displaystyle
 \frac{K_{m+2}}{K_m}
   =\frac{2^{8m+9}(m+1)\pi^{8m+7}}
              {(2m+2)!(2m+4)!}.}
\end{equation}
Each $K_m$ is a positive rational multiple of
$\pi^{2m^2-\lceil m/2\rceil}$.
\end{theorem}

\begin{proof}
Expand $\det T_m$ along the last row, using
\eqref{gn:eq:source-top}.  Since
$\phi_r(x)=rx^{2r-1}$ plus lower odd powers, triangular row elimination
in the remaining $(m-1)\times(m-1)$ minor extracts
\[
                  \prod_{r=1}^{m-1}\frac{4r}{(2r)!}.
\]
Reversing the source sinh-column order $m-1,\ldots,1$ to the increasing
order $1,\ldots,m-1$ contributes the sign
$(-1)^{(m-1)(m-2)/2}$.  The remaining determinant is exactly
$\mathcal D_{m-1}(1/4,\ldots,1/4)$, so
\begin{equation}\label{gn:eq:detT-factorization}
\begin{split}
 \det T_m={}&(-1)^{m+1+(m-1)(m-2)/2}
 \frac{2^{2m+2}\pi^m}{(2m)!}
   \left(\prod_{r=1}^{m-1}\frac{4r}{(2r)!}\right)\\
 &\hspace{24mm}\times\mathcal D_{m-1}(1/4,\ldots,1/4).
\end{split}
\end{equation}
The sign equals $(-1)^{m(m-1)/2}$: the exponents differ by two.
Theorem~\ref{gn:thm:determinant}, with $2m$ quarter parameters, gives
\begin{equation}\label{gn:eq:quarter-minor}
 \mathcal D_{m-1}(1/4,\ldots,1/4)
 =2^{2(m-1)^2}\pi^{m(m-1)}
       \frac{\Gamma(1/2)^{m(2m-1)}}{\Gamma(m/2)}.
\end{equation}
Substitute $\Gamma(1/2)=\sqrt\pi$, compare with
\eqref{gn:eq:source-normalization}, and collect the factorials and
powers.  This yields \eqref{gn:eq:K}.  At $m=1$ the minor is empty,
so the same computation remains valid.  Taking the quotient of the
formulas at $m+2$ and $m$, and using
$\Gamma(m/2+1)=(m/2)\Gamma(m/2)$, proves
\eqref{gn:eq:K-recurrence}.  The final assertion follows from the
integer and half-integer Gamma evaluations.
\end{proof}

The result makes the inherited exact Gamma determinant representation
fully explicit.  For example, if $F_{b,m}$ is the rational fundamental
matrix of the source recurrence, its already proved identity becomes
\begin{equation}\label{gn:eq:F-explicit}
 \det F_{b,m}
 =(-1)^{m(m-1)/2}
   \frac{\Gamma(m/2)\prod_{r=1}^m(2r)!}
        {2^{2m^2+2}(m-1)!\pi^{m(4m-1)/2}}
       \det(t_{b+i,j})_{i,j=1}^m.
\end{equation}
This evaluates a previously unspecified normalization.  It preserves
the source's distinction between its eventual fixed-gap nonvanishing
theorem and its still open all-pairs reconstruction conjecture.


% END INCLUDED SOURCE: sections/01_gamma

% BEGIN INCLUDED SOURCE: sections/02_gaussian
\section{Paired Gamma deformations and polygamma integrals}
\label{gn:sec:paired}

The parameter freedom in Theorem~\ref{gn:thm:determinant} is useful
beyond the single quarter specialization.  For $n\ge1$, $\sigma>0$,
and $m=n+1$, take
\begin{equation}\label{gn:eq:paired-parameters}
 a_{2r-1}=\sigma+u_r,\qquad a_{2r}=\sigma-u_r,
 \qquad 1\le r\le m,\qquad |\Re u_r|<\sigma.
\end{equation}
Let $\mathcal D_{n,\sigma}(\mathbf u)$ denote the corresponding
determinant.  Every pair has the same sum $2\sigma$, so the denominator
in \eqref{gn:eq:determinant} remains fixed.

\begin{theorem}[Complete paired deformation]
\label{gn:thm:paired}
On the domain \eqref{gn:eq:paired-parameters},
\begin{equation}\label{gn:eq:paired-gamma}
 \boxed{\displaystyle
 \frac{\mathcal D_{n,\sigma}(\mathbf u)}
      {\mathcal D_{n,\sigma}(\mathbf0)}
 =\prod_{r<s}
   \frac{\displaystyle\prod_{\varepsilon,\eta\in\{-1,1\}}
          \Gamma(2\sigma+\varepsilon u_r+\eta u_s)}
        {\Gamma(2\sigma)^4}.}
\end{equation}
At the quarter parameter this is elementary:
\begin{equation}\label{gn:eq:paired-sec}
 \boxed{\displaystyle
 \frac{\mathcal D_{n,1/4}(\mathbf u)}
      {\mathcal D_{n,1/4}(\mathbf0)}
 =\prod_{r<s}\sec\bigl(\pi(u_r+u_s)\bigr)
                 \sec\bigl(\pi(u_r-u_s)\bigr).}
\end{equation}
If only one pair is deformed, the two respective ratios are
\begin{equation}\label{gn:eq:one-pair}
 \left\{\frac{\Gamma(2\sigma+u)\Gamma(2\sigma-u)}
                    {\Gamma(2\sigma)^2}\right\}^{2n},
 \qquad \sec^{2n}(\pi u).
\end{equation}
\end{theorem}

\begin{proof}
The $m$ within-pair Gamma factors are $\Gamma(2\sigma)$ and are
unchanged.  Each pair of distinct labels contributes the four factors
shown in \eqref{gn:eq:paired-gamma}; all remaining constants cancel.
At $\sigma=1/4$, apply
$\Gamma(1/2+z)\Gamma(1/2-z)=\pi/\cos(\pi z)$ first to
$z=u_r+u_s$ and then to $z=u_r-u_s$.  The single-pair specialization
has $n$ identical pair interactions and gives \eqref{gn:eq:one-pair}.
\end{proof}

\subsection{A positive density and its exact normalization}

Define
\begin{equation}\label{gn:eq:density}
\begin{split}
 \rho_{n,\sigma}(\mathbf x)
 ={}&\Delta(x_1^2,\ldots,x_n^2)
      \Delta(\cosh\pi x_1,\ldots,\cosh\pi x_n)\\
 &\times\prod_{j=1}^n x_j\sinh(\pi x_j)
              |\Gamma(\sigma+ix_j/2)|^{4(n+1)},
 \qquad \mathbf x\in(0,\infty)^n.
\end{split}
\end{equation}
The density is symmetric and nonnegative, and is positive off the
collision hyperplanes.  Its normalizing integral is
\begin{equation}\label{gn:eq:Z}
 \mathcal Z_{n,\sigma}
 =2^{n(3n+1)/2}\pi^{n(n+1)}n!
       \frac{\Gamma(2\sigma)^{(n+1)(2n+1)}}
            {\Gamma(2(n+1)\sigma)}.
\end{equation}
This follows from \eqref{gn:eq:I-evaluation}.  Write $\mathbb E_{n,\sigma}$
for expectation with respect to $\rho_{n,\sigma}/\mathcal Z_{n,\sigma}$,
and define the real linear statistics
\begin{equation}\label{gn:eq:statistics}
 S_q=\sum_{j=1}^n\Re\psi^{(q)}(\sigma+ix_j/2),
                   \qquad q=1,3,5,\ldots.
\end{equation}

\begin{theorem}[All-order polygamma insertion identity]
\label{gn:thm:polygamma-germ}
As an identity of holomorphic germs at $\mathbf u=\mathbf0$,
\begin{equation}\label{gn:eq:polygamma-germ}
\begin{split}
 &\mathbb E_{n,\sigma}
  \exp\left(4\sum_{r=1}^{n+1}\sum_{k\ge1}
                   \frac{u_r^{2k}}{(2k)!}S_{2k-1}\right)\\
 &\hspace{12mm}=
 \prod_{r<s}
   \frac{\displaystyle\prod_{\varepsilon,\eta\in\{-1,1\}}
         \Gamma(2\sigma+\varepsilon u_r+\eta u_s)}
        {\Gamma(2\sigma)^4}.
\end{split}
\end{equation}
In particular, every finite coefficient gives an ordinary absolutely
convergent integral identity in polygamma functions.
\end{theorem}

\begin{proof}
Apply continuous Cauchy--Binet to the deformed determinant and divide
by its value at zero.  The interaction Vandermondes are unchanged, so
the quotient is the expectation of the product of Gamma-factor ratios.
For one pair and $z=\sigma+ix/2$, the logarithm of that ratio is
\[
\begin{aligned}
 &\log\Gamma(z+u)+\log\Gamma(z-u)-2\log\Gamma(z)\\
 &\quad+\log\Gamma(\overline z+u)+\log\Gamma(\overline z-u)
                         -2\log\Gamma(\overline z)\\
 &\hspace{18mm}=4\sum_{k\ge1}\frac{u^{2k}}{(2k)!}
                                      \Re\psi^{(2k-1)}(z).
\end{aligned}
\]
The logarithmic Gamma branches are holomorphic on the right half-plane.
Their Taylor expansions converge uniformly for $|u|$ in a fixed
sufficiently small disk, with the integral dominated as in
Theorem~\ref{gn:thm:determinant}.  Alternatively every fixed coefficient
can be obtained directly by differentiating the finite Gamma product
under its normally convergent integral.  Multiplication over spatial
variables and pair labels gives the left side of
\eqref{gn:eq:polygamma-germ}; Theorem~\ref{gn:thm:paired} evaluates it.
\end{proof}

For instance, the single-pair coefficients at degrees two and four are
\begin{align}
 \mathbb E_{n,\sigma}S_1
      &=n\psi'(2\sigma),\label{gn:eq:S1-mean}\\
 \mathbb E_{n,\sigma}(S_3+12S_1^2)
      &=n\psi'''(2\sigma)+12n^2\psi'(2\sigma)^2.
          \label{gn:eq:fourth-insertion}
\end{align}
The complete paired deformation separates the two terms on the left
of the second identity and yields a stronger moment law.

\section{Exact finite Gaussian moments}
\label{gn:sec:gaussian}

\begin{theorem}[Gaussian moment matching through degree $n+1$]
\label{gn:thm:gaussian}
Let $n\ge1$, $\sigma>0$, and set
\begin{equation}\label{gn:eq:mu-v}
 \mu_\sigma=\psi'(2\sigma)>0,\qquad
 v_\sigma=\psi'''(2\sigma)>0,
 \qquad S=S_1.
\end{equation}
For every integer $0\le k\le n+1$,
\begin{equation}\label{gn:eq:S-moments}
 \boxed{\displaystyle
 \mathbb E_{n,\sigma}S^k
 =\sum_{j=0}^{\lfloor k/2\rfloor}
       \frac{k!}{2^jj!(k-2j)!}
       (n\mu_\sigma)^{k-2j}(v_\sigma/4)^j.}
\end{equation}
Equivalently, the statistic
\begin{equation}\label{gn:eq:X}
             X_{n,\sigma}
       =\frac{2S-2n\mu_\sigma}{\sqrt{v_\sigma}}
\end{equation}
has exactly the standard Gaussian moments through degree $n+1$:
\begin{equation}\label{gn:eq:Gaussian-moments}
 \mathbb E_{n,\sigma}X_{n,\sigma}^{k}
 =\begin{cases}
    0,&k\text{ odd},\\[2pt]
    (k-1)!!,&k\text{ even},
   \end{cases}
 \qquad 0\le k\le n+1,
\end{equation}
where the moment at degree zero is one.
\end{theorem}

\begin{proof}
Choose $k$ distinct pair labels in \eqref{gn:eq:polygamma-germ} and
extract the coefficient of $u_1^2\cdots u_k^2$.  Put $w_r=u_r^2$ and
compute in the quotient in which $w_r^2=0$.  On the left, all terms
with $u_r^4$ or higher disappear, so the required coefficient is
$\mathbb E_{n,\sigma}(2S)^k$.

For one pair interaction, Taylor expansion of its logarithm gives
\begin{align*}
 &\sum_{\varepsilon,\eta\in\{-1,1\}}
      \log\Gamma(2\sigma+\varepsilon u_r+\eta u_s)
                       -4\log\Gamma(2\sigma)\\
 &\hspace{12mm}\equiv
          2\mu_\sigma(w_r+w_s)+v_\sigma w_rw_s
                              \pmod{w_r^2,w_s^2}.
\end{align*}
The sum is even in each pair variable.  Each of the $n+1$ labels has
exactly $n$ partners, and no logarithmic summand depends on three
distinct labels.  Thus the logarithm of the complete right side reduces
to
\[
       2n\mu_\sigma\sum_{r=1}^kw_r
                  +v_\sigma\sum_{r<s\le k}w_rw_s.
\]
To obtain $w_1\cdots w_k$ from its exponential, partition the $k$
labels into singleton and pair blocks.  The number of partitions with
$j$ pairs is $k!/(2^jj!(k-2j)!)$.  Such a partition contributes
$(2n\mu_\sigma)^{k-2j}v_\sigma^j$.  Therefore
\[
 \mathbb E_{n,\sigma}(2S)^k
 =\sum_{j=0}^{\lfloor k/2\rfloor}
   \frac{k!}{2^jj!(k-2j)!}
       (2n\mu_\sigma)^{k-2j}v_\sigma^j.
\]
Division by $2^k$ proves \eqref{gn:eq:S-moments}.  These are the
moments of a Gaussian variable with mean $n\mu_\sigma$ and variance
$v_\sigma/4$.  Centering and scaling proves
\eqref{gn:eq:Gaussian-moments}.  The restriction $k\le n+1$ is exactly
the number of distinct pair parameters available in the proof.
\end{proof}

\begin{corollary}[Separated low-order polygamma integrals]
\label{gn:cor:variance}
For all $n\ge1$ and $\sigma>0$,
\begin{equation}\label{gn:eq:variance}
 \operatorname{Var}_{n,\sigma}(S_1)=\frac14\psi'''(2\sigma),
 \qquad
 \mathbb E_{n,\sigma}S_3=(n-3)\psi'''(2\sigma).
\end{equation}
At the quarter parameter,
\begin{equation}\label{gn:eq:quarter-statistics}
 \mathbb E_{n,1/4}S_1=\frac{n\pi^2}{2},\qquad
 \operatorname{Var}_{n,1/4}(S_1)=\frac{\pi^4}{4},\qquad
 \mathbb E_{n,1/4}S_3=(n-3)\pi^4.
\end{equation}
\end{corollary}

\begin{proof}
The degree-two instance of \eqref{gn:eq:S-moments} gives the variance.
Insert the resulting second moment into
\eqref{gn:eq:fourth-insertion} to obtain the expectation of $S_3$.
Finally, $\psi'(1/2)=\pi^2/2$ and $\psi'''(1/2)=\pi^4$ give the
quarter values \cite{dlmfpolygamma}.
\end{proof}

\subsection{Explicit one-dimensional identities}

At $n=1$ and $\sigma=1/4$, the normalizer is $4\pi^5$.  To display
several resulting evaluations compactly, write
\[
 \omega(x)=x\sinh(\pi x)|\Gamma(\tfrac14+ix/2)|^8,
 \qquad z_x=\tfrac14+ix/2.
\]
Every integral below is an ordinary absolutely convergent integral:
\begin{align}
 \int_0^\infty\omega(x)\dd x&=4\pi^5,
          \label{gn:eq:scalar-normalizer}\\
 \int_0^\infty\omega(x)\Re\psi'(z_x)\dd x&=2\pi^7,
          \label{gn:eq:scalar-trigamma}\\
 \int_0^\infty\omega(x)[\Re\psi'(z_x)]^2\dd x&=2\pi^9,
          \label{gn:eq:scalar-trigamma-square}\\
 \int_0^\infty\omega(x)\Re\psi'''(z_x)\dd x&=-8\pi^9.
          \label{gn:eq:scalar-third}
\end{align}
The negative sign in \eqref{gn:eq:scalar-third} is compatible with an
ordinary real integral: the real part of the complex-argument
polygamma function changes sign.

Unpaired parameter derivatives supply digamma identities as well.
Differentiating the $n=1$ version of
\eqref{gn:eq:determinant} once in one parameter and twice in two
distinct parameters, then setting all four parameters equal to $1/4$,
gives
\begin{align}
 \int_0^\infty\omega(x)\Re\psi(z_x)\dd x
    &=-4\pi^5(\gamma+3\log2),\label{gn:eq:scalar-psi}\\
 \int_0^\infty\omega(x)[\Re\psi(z_x)]^2\dd x
    &=4\pi^5\left((\gamma+3\log2)^2+\frac{\pi^2}{12}\right).
       \label{gn:eq:scalar-psi-square}
\end{align}
For completeness, the first logarithmic parameter derivative of the
Gamma quotient is $3\psi(1/2)-\psi(1)$, while the mixed logarithmic
derivative in distinct labels is $\psi'(1/2)-\psi'(1)$.  The integral
insertions are respectively $2\Re\psi(z_x)$ and
$4[\Re\psi(z_x)]^2$.  These factors give exactly
\eqref{gn:eq:scalar-psi}--\eqref{gn:eq:scalar-psi-square}.

\subsection{A distributional consequence of the exact identities}

The Gaussian matching is asserted only through the stated finite degree.
For fixed $n$ and $\sigma$ the statistic is bounded, since
\[
 |\Re\psi'(\sigma+iy)|\le\sum_{r=0}^\infty\frac1{(r+\sigma)^2}
                       =\psi'(\sigma).
\]
Thus it has a compactly supported distribution.  Its moment identities
nevertheless imply a Gaussian limit as the dimension increases.

\begin{corollary}[Gaussian limit for arbitrary positive parameter sequences]
\label{gn:cor:CLT}
For any sequence $\sigma_n>0$, the random variables
$X_{n,\sigma_n}$ of \eqref{gn:eq:X} converge in distribution to a
standard normal random variable as $n\to\infty$.
\end{corollary}

\begin{proof}
For every fixed integer $k$, the $k$th moment is exactly the standard
Gaussian moment once $n\ge k-1$, independently of $\sigma_n$.
The second moment equals one for every $n\ge1$, giving tightness.
For each fixed $k$, a higher even moment is uniformly bounded along
the tail of the sequence, so the $k$th powers are uniformly integrable.
Every subsequential weak limit therefore has all standard Gaussian
moments.  The normal distribution is moment determinate; for example
its even moments satisfy Carleman's condition.  Consequently every
subsequential limit is standard normal, proving the assertion.
\end{proof}

This consequence uses exact moment identities at each sufficiently
large finite dimension.  No uniform large-dimension estimate on Gamma
functions or on the parameter sequence is required.

% END INCLUDED SOURCE: sections/02_gaussian

% BEGIN INCLUDED SOURCE: sections/03_colored
\section{Colored balanced kernels and their Mellin transform}
\label{sec:colored}

Let
\[
 \Omega=\C\setminus[1,\infty),\qquad a>0,\qquad x>0,\qquad \omega\in\Omega.
\]
The branch of $\Phi(\omega,s,a)$ is determined by its defining series
for $|\omega|<1$ and continuation to $\Omega$. Equivalently, for
$\Re s>0$ it is given by the classical absolutely convergent integral
\cite[25.14.5]{dlmflerch}
\begin{equation}\label{col:lerch-integral}
 \Phi(\omega,s,a)=\frac1{\Gamma(s)}
  \int_0^\infty\frac{t^{s-1}e^{-at}}{1-\omega e^{-t}}\dd t.
\end{equation}
Subtracting Taylor coefficients of its smooth density at zero extends
it to an entire function of $s$ for each $\omega\in\Omega$.

\begin{definition}[Colored balanced derivative]\label{col:def}
For an integer $q\ge0$, put
\begin{align}
 F_\omega(z;x,a)&=x^z\Phi(\omega,1+z,x+a),\label{col:generator}\\
 C_q^\omega(x,a)&=\left.\partial_z^qF_\omega(z;x,a)\right|_{z=0}\notag\\
 &=\sum_{j=0}^q\binom qj\log^{q-j}x\,
       \partial_s^j\Phi(\omega,1,x+a).\label{col:kernel}
\end{align}
Order derivatives are evaluated after analytic continuation in $\omega$.
\end{definition}

There is no order pole at $s=1$ when $\omega\ne1$ in this domain.
Thus adding the uncolored term $\log^{q+1}x/(q+1)$ to
\eqref{col:kernel} would destroy decay rather than improve it.

\subsection{Tail cancellation and uniqueness}

Write
\begin{equation}\label{col:A}
 A_k(\omega,a)=\Phi(\omega,-k,a)
       =(\omega\partial_\omega+a)^k\frac1{1-\omega}.
\end{equation}
These are rational functions of $\omega$ and polynomials of degree $k$
in $a$. Their leading coefficient in $a$ is $1/(1-\omega)$.

\begin{lemma}[Uniform colored expansion]\label{col:tail}
On compact subsets of $\Omega\times(0,\infty)$, and locally uniformly
for $z$ near zero,
\begin{equation}\label{col:F-tail}
 F_\omega(z;x,a)=
  \sum_{k=0}^{K}\frac{(-1)^k A_k(\omega,a)}{k!}
           (1+z)_k x^{-k-1}+O(x^{-K-2}).
\end{equation}
Consequently
\begin{equation}\label{col:C-tail}
 C_q^\omega(x,a)=
  \sum_{k=q}^{K}\frac{(-1)^k A_k(\omega,a)}{k!}
   \left.\partial_z^q(1+z)_k\right|_{z=0}x^{-k-1}
       +O(x^{-K-2}),
\end{equation}
and the leading coefficient is
$(-1)^q A_q(\omega,a)x^{-q-1}$.
At zero, $C_q^\omega(x,a)=O(1+|\log x|^q)$.
\end{lemma}

\begin{proof}
The density $e^{-at}/(1-\omega e^{-t})$ is analytic at $t=0$,
with Taylor coefficients $(-1)^k A_k(\omega,a)/k!$.
Taylor's formula on a small fixed interval and exponential decay
outside that interval give a Laplace expansion in
\eqref{col:lerch-integral}. The $k$th term integrates to
$\Gamma(1+z+k)x^{-1-z-k}$; division by $\Gamma(1+z)$
and multiplication by $x^z$ give \eqref{col:F-tail}.
On the stated compact sets, the denominator stays uniformly away
from zero along $t\ge0$, so the remainders are locally uniform.
Cauchy differentiation in a slightly larger $z$ disc preserves
the inverse-power order of the remainder. Since $(1+z)_k$ has
degree $k$, derivatives of order $q>k$ vanish, and its leading
coefficient gives $\partial_z^q(1+z)_q=q!$.
The assertion at zero follows directly from \eqref{col:kernel}.
\end{proof}

\begin{theorem}[Unique homogeneous balancing]\label{col:unique}
Fix $\omega\in\Omega$ and $q\ge0$. Among combinations
\[
 D(x,a)=\sum_{j=0}^q d_j\log^{q-j}x\,
                  \partial_s^j\Phi(\omega,1,x+a),\qquad d_0=1,
\]
with coefficients independent of $x,a$, the condition
$D(x,a)=O(x^{-q-1})$ for every $a$ in a nonempty open interval contained in $(0,\infty)$
forces $d_j=\binom qj$. These coefficients do satisfy that condition.
\end{theorem}

\begin{proof}
Put $d(t)=\sum_{j=0}^q d_jt^j$ and $L=\log x$.
The coefficient arising from the $k$th density term is a polynomial
in $L$. Its component with $L^{q-r}$ contains
$d^{(r)}(-1)$ times a nonzero scalar and
$\partial_s^r(s)_k|_{s=1}$. For $k=r<q$, the latter derivative
is $r!$. The polynomial $A_r(\omega,a)$ cannot vanish on the
whole prescribed interval. Successively removing $x^{-1},\ldots,x^{-q}$
therefore gives $d^{(r)}(-1)=0$ for $r=0,\ldots,q-1$.
Since $\deg d\le q$ and $d(0)=1$, $d(t)=(1+t)^q$.
The converse is Lemma~\ref{col:tail}. This proves uniqueness only
in the displayed ansatz, not among arbitrary $x$-dependent combinations.
\end{proof}

% Optional insertion after Theorem col:unique. Uses root's notation.
\begin{remark}[The interval in the shift is essential]
\label{col:pointwise-counterexample}
The uniqueness statement can fail if imposed only at one shift.
Set
\[
 a_* = \frac{1+\sqrt2}{2},\qquad
 \omega_*=-3-2\sqrt2\in\Omega.
\]
The rational second moment factors as
\[
 A_2(\omega,a_*)
  =\frac{(3-2\sqrt2)(\omega+3+2\sqrt2)^2}
          {4(1-\omega)^3}.
\]
Consequently $A_2(\omega_*,a_*)=0$ and, by
$A_3=(\omega\partial_\omega+a_*)A_2$,
$A_3(\omega_*,a_*)=0$ as well.
The uniform tail formula therefore gives
\[
 C_2^{\omega_*}(x,a_*)=O(x^{-5}),\qquad
 C_3^{\omega_*}(x,a_*)=O(x^{-5}).
\]
Hence the different homogeneous combination
\[
 D(x)=2C_3^{\omega_*}(x,a_*)
                  -\log x\,C_2^{\omega_*}(x,a_*)
\]
also satisfies $D(x)=O(x^{-4})$.  In the displayed uniqueness
ansatz its coefficient vector is
$\boldsymbol d=(1,4,5,2)$ rather than $(1,3,3,1)$, so it has
the same normalization $d_0=1$.  The fourth moment is nonzero,
$A_4(\omega_*,a_*)=(2-\sqrt2)/64$, and the modified kernel is
not identically the canonical one.  Requiring decay for a
nonempty open interval of positive shifts excludes this
pointwise degeneracy exactly as the proof of the theorem requires.
\end{remark}

\subsection{The exact Mellin identity}

Define the generalized harmonic polynomial $Q_q(u)$ by
\begin{equation}\label{col:Q}
 \frac{\Gamma(u+z)}{\Gamma(u)\Gamma(1+z)}
     =\sum_{q=0}^\infty Q_q(u)\frac{z^q}{q!}.
\end{equation}
With
$H_{u-1}^{(1)}=\psi(u)+\gamma$ and
$H_{u-1}^{(r)}=\zeta(r)-\zeta(r,u)$ for $r\ge2$, this is
\[
 \exp\!\left(\sum_{r=1}^\infty
          \frac{(-1)^{r-1}}r H_{u-1}^{(r)}z^r\right).
\]
In particular,
$Q_0=1$, $Q_1=H^{(1)}$, and
$Q_2=(H^{(1)})^2-H^{(2)}$.

\begin{theorem}[Colored balanced Mellin formula]\label{col:mellin}
For $q\ge0$ and $0<\Re u<q+1$,
\begin{equation}\label{col:mellin-formula}
 \boxed{\displaystyle
 \int_0^\infty x^{u-1}C_q^\omega(x,a)\dd x
  =\frac{\pi}{\sin\pi u}\,Q_q(u)\Phi(\omega,1-u,a).}
\end{equation}
Every apparent singularity at an interior integer is removable.
The integral is absolutely convergent and locally uniform in
$u,\omega,a$ on the stated domains.
\end{theorem}

\begin{proof}
Initially take $0<\Re u<1$ and $z$ sufficiently close to zero.
Insert \eqref{col:lerch-integral} into \eqref{col:generator}.
Absolute Fubini and the Gamma integral give
\begin{align*}
 \int_0^\infty x^{u-1}F_\omega(z;x,a)\dd x
 &=\frac{\Gamma(u+z)}{\Gamma(1+z)}
       \int_0^\infty\frac{t^{-u}e^{-at}}{1-\omega e^{-t}}\dd t\\
 &=\frac{\Gamma(u+z)\Gamma(1-u)}{\Gamma(1+z)}
                  \Phi(\omega,1-u,a).
\end{align*}
The behavior at zero and infinity in Lemma~\ref{col:tail}
justifies locally normal differentiation in $z$.
Euler reflection then proves \eqref{col:mellin-formula}
on the initial strip. The differentiated kernel decays as
$x^{-q-1}$, so its integral is holomorphic on the larger
strip $0<\Re u<q+1$. Analytic continuation establishes the
formula there and makes the interior integer poles removable.
This does not differentiate a divergent integral at an integer.
\end{proof}

\section{Integer colored moments, Stieltjes differences, and primitives}
\label{sec:colored-integer}

For $q\ge m\ge1$, let
\begin{align}
 L_{m,q}
 &:=(-1)^q q!\sum_{\ell=m}^\infty
              \frac{\binom{\ell-1}{m-1}}{\ell^{q+1}}
 \label{col:L-positive}\\
 &=(-1)^{m+q+1}q!
       \sum_{j=0}^{m-1}(-1)^j e_j(1,\tfrac12,\ldots,\tfrac1{m-1})
                         \zeta(q+1-j).\label{col:L-finite}
\end{align}
The first representation is convergent and gives its sign.
The second is a finite zeta expression, with the empty-list convention
$e_0=1$.

\begin{theorem}[All ordinary positive-integer moments]\label{col:integer}
For $q\ge m\ge1$,
\begin{equation}\label{col:integer-formula}
 \boxed{\displaystyle
 \int_0^\infty x^{m-1}C_q^\omega(x,a)\dd x
             =L_{m,q}\Phi(\omega,1-m,a).}
\end{equation}
Thus the complete shift dependence is a polynomial of degree
$m-1$, with rational functions of the color as coefficients.
\end{theorem}

\begin{proof}
At $u=m$, the Gamma quotient in \eqref{col:Q} is
$\prod_{j=1}^{m-1}(1+z/j)$, hence $Q_q(m)=0$ for $q\ge m$.
The limit in \eqref{col:mellin-formula} is
$(-1)^m Q_q'(m)\Phi(\omega,1-m,a)$.
Differentiating the quotient in $u$ gives
\[
 Q_q'(m)=q![z^q]\prod_{j=1}^{m-1}(1+z/j)
                      \bigl(\psi(m+z)-\psi(m)\bigr).
\]
The rational terms in the digamma recurrence become polynomials
of degree at most $m-1$ after multiplication by the finite product.
They do not contribute at degree $q$. Expanding
$\psi(1+z)-\psi(1)=\sum_{r\ge1}(-1)^{r+1}\zeta(r+1)z^r$
proves $(-1)^mQ_q'(m)=L_{m,q}$ in the finite form.
Expanding the zeta values and using
\[
 \prod_{j=1}^{m-1}(1-\ell/j)
       =(-1)^{m-1}\binom{\ell-1}{m-1}
\]
gives \eqref{col:L-positive}.
\end{proof}

There is an independent elementary proof for $|\omega|<1$.
In that disk,
\begin{equation}\label{col:elementary-series}
 C_q^\omega(x,a)=\sum_{\ell=0}^\infty
  \frac{\omega^\ell}{x+\ell+a}
           \log^q\!\frac{x}{x+\ell+a}.
\end{equation}
For $c>0$, the substitution $t=x/(x+c)$ gives
\begin{equation}\label{col:elementary-moment}
 \int_0^\infty\frac{x^{m-1}}{x+c}
                  \log^q\!\frac{x}{x+c}\dd x
 =c^{m-1}\int_0^1\frac{t^{m-1}\log^q t}{(1-t)^m}\dd t
 =c^{m-1}L_{m,q}.
\end{equation}
Absolute summation proves the theorem in the disk, and locally normal
continuation in $\omega$ gives the rest of $\Omega$. At $\omega=-1$,
one must not replace that continuation by the divergent series
$\sum(-1)^\ell(\ell+a)^{m-1}$.

\subsection{Low-order examples and logarithmic insertions}

For every $q\ge1$ the unweighted moment is independent of the shift:
\begin{equation}\label{col:unweighted}
 \boxed{\displaystyle
 \int_0^\infty C_q^\omega(x,a)\dd x
       =\frac{(-1)^q q!\zeta(q+1)}{1-\omega}.}
\end{equation}
For example,
\begin{align}
 \int_0^\infty C_1^\omega(x,a)\dd x
      &=-\frac{\zeta(2)}{1-\omega},\label{col:first}\\
 \int_0^\infty xC_2^\omega(x,a)\dd x
      &=2\bigl(\zeta(2)-\zeta(3)\bigr)
         \left(\frac{a}{1-\omega}
                    +\frac{\omega}{(1-\omega)^2}\right).
 \label{col:second}
\end{align}
All fixed logarithmic insertions are equally explicit. For $q\ge m$
and $r\ge0$,
\begin{equation}\label{col:log-moments}
\begin{split}
 &\int_0^\infty x^{m-1}\log^r x\,C_q^\omega(x,a)\dd x\\
 &\quad=(-1)^m r![t^r]\left\{
   \frac{\pi t}{\sin\pi t}\,
   \frac{Q_q(m+t)}t\,\Phi(\omega,1-m-t,a)\right\}.
\end{split}
\end{equation}
The quotient $Q_q(m+t)/t$ is holomorphic. This is a finite coefficient
formula in Lerch order derivatives, Gamma derivatives, and generalized
harmonic polynomials. It follows by differentiating the normally
convergent Mellin integral, not by termwise integration of its
individual divergent logarithmic summands.

\subsection{A precise generalized Stieltjes dictionary}

Let $\omega^d=1$ with $\omega\ne1$. Introduce the incoming uncolored
kernel \cite{repobalanced}
\begin{equation}\label{col:uncolored}
 C_q(x,a)=\frac{\log^{q+1}x}{q+1}
       +\sum_{j=0}^q(-1)^j\binom qj
                  \gamma_j(x+a)\log^{q-j}x.
\end{equation}
The finite residue-class identity for Lerch's function gives
\begin{equation}\label{col:stieltjes-dictionary}
 \boxed{\displaystyle
 C_q^\omega(x,a)=\frac1d\sum_{r=0}^{d-1}
       \omega^r C_q\!\left(\frac xd,\frac{a+r}{d}\right).}
\end{equation}
Indeed
$\Phi(\omega,s,x+a)=d^{-s}\sum_r\omega^r
\zeta(s,(x+a+r)/d)$; the $1/z$ poles and the
$\log^{q+1}(x/d)$ terms cancel because $\sum_r\omega^r=0$.
This identity translates every preceding colored formula into an
ordinary integral of a finite combination of generalized Stieltjes
constants. It also verifies the normalization against the incoming
finite-filter calculus.

For the alternating color, define
\[
 a_-=\frac{x+a}{2},\qquad a_+=\frac{x+a+1}{2}.
\]
Then, explicitly,
\begin{equation}\label{col:alternating-C1}
 C_1^{-1}(x,a)=\frac12\left[
    \bigl(\psi(a_+)-\psi(a_-)\bigr)\log\frac x2
                 +\gamma_1(a_+)-\gamma_1(a_-)\right].
\end{equation}
Its ordinary integral is $-\pi^2/12$ for every $a>0$.

\subsection{Exceptional moments and Dirichlet beta values}

\begin{theorem}[A zero of the leading colored tail]\label{col:exceptional}
If $m\ge1$ and $\Phi(\omega,1-m,a)=0$, then
\begin{equation}\label{col:exceptional-formula}
 \boxed{\displaystyle
 \int_0^\infty x^{m-1}C_{m-1}^\omega(x,a)\dd x
       =(-1)^{m+1}\partial_s\Phi(\omega,1-m,a).}
\end{equation}
The integral is ordinary and absolutely convergent.
\end{theorem}

\begin{proof}
The zero hypothesis removes the leading $x^{-m}$ kernel term in
Lemma~\ref{col:tail}, so the Mellin integral extends
holomorphically across $u=m$. At $u=m$, $Q_{m-1}(m)=1$ and
$\Phi(\omega,1-u,a)=-(u-m)\partial_s\Phi(\omega,1-m,a)
 +O((u-m)^2)$. The removable limit in
Theorem~\ref{col:mellin} is \eqref{col:exceptional-formula}.
\end{proof}

Let $\beta(s)=\sum_{n\ge0}(-1)^n(2n+1)^{-s}$ denote Dirichlet's beta
function, continued in its order. Since
$\Phi(-1,s,1/2)=2^s\beta(s)$, its negative odd zeros give the following
family:
\begin{corollary}[Alternating centered exceptional moments]\label{col:beta}
For every integer $r\ge1$,
\begin{equation}\label{col:beta-formula}
 \boxed{\displaystyle
 \int_0^\infty x^{2r-1}C_{2r-1}^{-1}(x,\tfrac12)\dd x
   =\frac{(-1)^r(2r-1)!}{\pi^{2r-1}}\beta(2r).}
\end{equation}
In particular, writing $G=\beta(2)$ for Catalan's constant,
\begin{equation}\label{col:catalan}
 \int_0^\infty xC_1^{-1}(x,\tfrac12)\dd x=-\frac G\pi.
\end{equation}
\end{corollary}

\begin{proof}
The completed beta function
$(4/\pi)^{(s+1)/2}\Gamma((s+1)/2)\beta(s)$ is invariant under
$s\mapsto1-s$ \cite{dlmfdirichlet}. Comparing its finite values at
$s=1-2r$ and $s=2r$ yields
\[
 \beta'(1-2r)=
 (-1)^{r-1}2^{2r-1}(2r-1)!\pi^{1-2r}\beta(2r).
\]
Substitute this into Theorem~\ref{col:exceptional}.
The classical functional equation is the input; the ordinary
Stieltjes-difference integral is its consequence in the new kernel.
\end{proof}

\subsection{Anchored shift primitives}

For $a,a_0>0$ the quotient below is regular at $z=0$, because
$\Phi(\omega,0,A)=1/(1-\omega)$ does not depend on $A$:
\begin{equation}\label{col:primitive}
\begin{split}
 \int_{a_0}^a C_q^\omega(x,v)\dd v
  =-q![z^q]\,\frac{x^z}{z}
       \left(\Phi(\omega,z,x+a)-\Phi(\omega,z,x+a_0)\right).
\end{split}
\end{equation}
This follows from
$\partial_A\Phi(\omega,z,A)=-z\Phi(\omega,z+1,A)$.
It gives every integration constant through the explicit anchor.
At order zero and a root of unity, the Hurwitz identity
$\zeta'(0,A)=\log\Gamma(A)-\tfrac12\log(2\pi)$ gives the finite
Gamma expression
\begin{equation}\label{col:gamma-primitive}
 \int_{a_0}^a C_0^\omega(x,v)\dd v
  =-\sum_{r=0}^{d-1}\omega^r
       \log\frac{\Gamma((x+a+r)/d)}
                    {\Gamma((x+a_0+r)/d)}.
\end{equation}
All Gamma arguments are positive, and their logarithms are real
before multiplication by the complex coefficients.

\section{The color-one limit and its exact constant term}
\label{sec:color-limit}

The formulas above are not uniform at $\omega=1$. The precise
normalization can nevertheless be determined, and it reconciles the
colored moment with its uncolored Stieltjes counterpart.
Fix $q\ge m\ge1$, let $\omega=e^{-\lambda}$ with $\lambda>0$, and set
\[
 M_{m,q}(\lambda,a)=\int_0^\infty
                   x^{m-1}C_q^{e^{-\lambda}}(x,a)\dd x.
\]

\begin{theorem}[Complete local color expansion]\label{col:coalescence}
For $0<\lambda<2\pi$,
\begin{equation}\label{col:coalescence-formula}
 \boxed{\displaystyle
 e^{-a\lambda}M_{m,q}(\lambda,a)
 =L_{m,q}\left\{\frac{\Gamma(m)}{\lambda^m}
       +\sum_{k=0}^\infty
             \frac{(-\lambda)^k}{k!}\zeta(1-m-k,a)\right\}.}
\end{equation}
The series is convergent. In the cutoff coordinate $\lambda$,
\begin{align}
 \lim_{\lambda\downarrow0}
 \left(e^{-a\lambda}M_{m,q}
              -L_{m,q}\Gamma(m)\lambda^{-m}\right)
      &=-\frac{L_{m,q}}m B_m(a),\label{col:renorm}\\
 \CT_{\lambda=0}M_{m,q}(\lambda,a)
      &=\frac{L_{m,q}}m\bigl(a^m-B_m(a)\bigr).\label{col:raw-CT}
\end{align}
The first line is exactly the uncolored ordinary moment
$\int_0^\infty x^{m-1}C_q(x,a)\dd x$ from \cite{repobalanced}.
\end{theorem}

\begin{proof}
The classical convergent Lerch expansion \cite[25.14.3\_1]{dlmflerch}
is
\[
 e^{-a\lambda}\Phi(e^{-\lambda},s,a)
    =\Gamma(1-s)\lambda^{s-1}
       +\sum_{k=0}^\infty\zeta(s-k,a)\frac{(-\lambda)^k}{k!}.
\]
Set $s=1-m$ and use Theorem~\ref{col:integer}.
Its constant term is $L_{m,q}\zeta(1-m,a)
=-L_{m,q}B_m(a)/m$.
Multiplying by $e^{a\lambda}$ adds the constant contribution
$L_{m,q}\Gamma(m)a^m/m!=L_{m,q}a^m/m$ from the pole.
This proves both normalizations.
The incoming uncolored coefficient is
$\kappa_{m,q}=-L_{m,q}/m$, so its moment is the first limit.
\end{proof}

For example, $M_{1,1}=-\zeta(2)/(1-e^{-\lambda})$ is independent of
$a$, and its raw constant term is $-\zeta(2)/2$.
The normalized limit in \eqref{col:renorm}, however, is
$\zeta(2)(a-\tfrac12)$. Both formulas are correct in their stated
normalizations. They cannot be identified by discarding the
$e^{-a\lambda}$ factor.

All logarithmic moment versions follow by differentiating the local
Mellin expression in $u$ before putting $u=m$. Derivatives of
$\lambda^{-u}$ produce explicit powers of $\log\lambda$.
Thus the scalar expansion already specifies the complete mechanism
for those contacts. It does not assert that the pointwise limit of
the colored kernels can be integrated across the entire half-line.


% END INCLUDED SOURCE: sections/03_colored

% BEGIN INCLUDED SOURCE: sections/04_clusters
% Requires amsmath, amssymb, amsthm and standard theorem environments.
% References to add: repolattice, okudaueno, dlmfbeta,
% dlmfgamma.  Bibitems and precise source coverage are in
% two_cluster_source_status.txt.  Every label has the prefix tc:.

\section{An explicit polylogarithm dictionary for two mass clusters}
\label{tc:sec:clusters}

The incoming \emph{Lattice Contacts and Mixed Zeta Identities} report
evaluates a Gamma face when all but one mass agree.  Its spectral-face
section then asks for a finite dictionary for two distinct mass
clusters; its research question~Q4 asks more broadly about displaced
masses and the dependence of depth on the number of clusters
\cite{repolattice}.  We evaluate the entire two-cluster sector,
with arbitrary positive cluster multiplicities.  The result is a
finite combination of ordinary one-variable multiple polylogarithms
inside an absolutely convergent spectral sum.  In particular, two
equal displaced masses require depth at most three, independently
of the total number of factors.

The argument uses classical beta integration, aggregation of Dirichlet
parameters, the Gamma product, and coefficient extraction.  The
transformation under exchange of clusters is a specialization of
classical multiple-polylogarithm Landen formulas
\cite{okudaueno}; it is not asserted to be a new general
transformation law.  The contribution relative to the inspected
incoming source is the explicit all-multiplicity dictionary and its
convergence and depth statement.  Minimality of the resulting depth
and a bound depending only on the number of distinct clusters remain
unproved.

\subsection{The ordinary face and its independent regulators}

Let $k\ge2$, $b_1,\ldots,b_k\in\mathbb R$, and
\[
 \Sigma_k=\{(t_1,\ldots,t_k):t_i\ge0,\ \sum_i t_i=1\}.
\]
The measure $dt$ means $dt_1\cdots dt_{k-1}$ after eliminating
$t_k$; it has no additional Euclidean surface factor.  For nonempty
$T\subseteq\{1,\ldots,k\}$ set
\[
 t_T=\sum_{i\in T}t_i,\qquad
 b_T=\frac{\sum_{i\in T}t_i b_i}{t_T},
\]
and give $t_TF(b_T)$ the value zero when $t_T=0$.  Define
\begin{align}
 \mathcal H_k(F;t)
  &=\sum_{\varnothing\ne T\subseteq\{1,\ldots,k\}}
       (-1)^{k-|T|}t_TF(b_T),\label{tc:eq:face-integrand}\\
 I_k(F;b_1,\ldots,b_k)
  &=\int_{\Sigma_k}\mathcal H_k(F;t)
                    \frac{dt}{t_1\cdots t_k}.
                    \label{tc:eq:ordinary-face}
\end{align}
Constants and affine functions are annihilated.

These are ordinary integrals.  To see convergence for $F\in C^1$
on the mass interval, each perspective is Lipschitz because
\[
 \partial_{t_i}\{t_TF(b_T)\}
    =F(b_T)+F'(b_T)(b_i-b_T),\qquad i\in T.
\]
The alternating sum vanishes on every coordinate face, so
$|\mathcal H_k(F;t)|\le C\min_i t_i$.
On the region where $t_j$ is maximal, $t_j\ge1/k$ and
\[
 \min_i t_i\le\left(\prod_{i\ne j}t_i\right)^{1/(k-1)}.
\]
Consequently the integral is absolutely convergent.  The same argument
allows sufficiently small independent complex powers $t_i^{\alpha_i}$
and any fixed powers of their logarithms.

For $\alpha_i>0$, $A=\sum_i\alpha_i$, let
\[
 D_F(\boldsymbol\alpha;\boldsymbol b)
   =\frac{\Gamma(A)}{\prod_i\Gamma(\alpha_i)}
      \int_{\Sigma_k}F\!\left(\sum_i t_i b_i\right)
                            \prod_i t_i^{\alpha_i-1}\,dt.
\]
If $A_T=\sum_{i\in T}\alpha_i$, the classical beta integral
\cite{dlmfbeta} gives
\begin{equation}\label{tc:eq:independent-regulators}
 \frac{\Gamma(1+A)}{\prod_i\Gamma(1+\alpha_i)}
 \int_{\Sigma_k}\mathcal H_k(F;t)
                     \prod_i t_i^{\alpha_i-1}\,dt
 =\frac{1}{\prod_i\alpha_i}
   \sum_{\varnothing\ne T}(-1)^{k-|T|}
                  A_TD_F(\boldsymbol\alpha_T;\boldsymbol b_T).
\end{equation}
Here the notation on the right denotes the Dirichlet mean formed
only from the labels in $T$, not the random barycenter on the left.
The identity follows from
$\mathbb E[t_TF(b_T)]=(A_T/A)D_F(\boldsymbol\alpha_T;\boldsymbol b_T)$.
It is important that the $\alpha_i$ are independent until the
squarefree coefficient is taken.  The left side is holomorphic near
zero and has value $I_k(F;\boldsymbol b)$ there.

\subsection{One scalar beta integral for arbitrary multiplicities}

Let $r,s\ge1$, $k=r+s$, and write $B^r,C^s$ for $r$ copies of
$B$ followed by $s$ copies of $C$.  Put $\delta=C-B$ and
\begin{align}
 E_F(t)&=F(B+\delta t)-(1-t)F(B)-tF(C),\label{tc:eq:chord}\\
 Q(X,Y)&=\frac{\Gamma(1+X+Y)}{\Gamma(1+X)\Gamma(1+Y)}.
                                                     \label{tc:eq:Q}
\end{align}

\begin{theorem}[Scalar beta representation for two clusters]
\label{tc:thm:beta}
For $F\in C^2$ on the interval with endpoints $B,C$,
\begin{equation}\label{tc:eq:beta-representation}
 \boxed{\displaystyle
 I_{r+s}(F;B^r,C^s)
   =r!s!\,[X^{r-1}Y^{s-1}]\,
      Q(X,Y)\int_0^1
      \frac{E_F(t)}{t(1-t)}(1-t)^X t^Y\,dt.}
\end{equation}
The integral on the right is holomorphic, for example for
$\Re X,\Re Y>-1$, and the coefficient may be taken under the integral.
\end{theorem}

\begin{proof}
Aggregate the independent regulators belonging to the $B$ cluster
into $X$ and those belonging to the $C$ cluster into $Y$.  The
proportion of the $C$ cluster has beta density
$t^{Y-1}(1-t)^{X-1}/\mathrm B(Y,X)$.  Thus the full-set term
$AD_F$ from \eqref{tc:eq:independent-regulators} becomes
\begin{align}
 \mathcal M_F(X,Y)
 &=\frac{\Gamma(1+X+Y)}{\Gamma(X)\Gamma(Y)}
     \int_0^1 F(B+\delta t)t^{Y-1}(1-t)^{X-1}\,dt\notag\\
 &=XF(B)+YF(C)
   +XYQ(X,Y)\int_0^1
          \frac{E_F(t)}{t(1-t)}(1-t)^Xt^Y\,dt.
                                                     \label{tc:eq:M}
\end{align}
The second line follows by subtracting the endpoint chord and
integrating it exactly.  Since $E_F(t)=O(t(1-t))$, it proves
holomorphy of the full-set term near $(0,0)$ without leaving
unremoved beta poles.

In the alternating sum in \eqref{tc:eq:independent-regulators},
each proper-subset term is holomorphic by the same chord
subtraction, or is linear when it contains only one cluster.
Every proper subset omits at least one regulator.  Its coefficient
of $\alpha_1\cdots\alpha_k$ is therefore zero.  In the full-set term,
the corresponding coefficient is
$r!s![X^rY^s]\mathcal M_F(X,Y)$: the squarefree coefficient in the
$r$th power of a sum of $r$ independent variables is $r!$.
The chord term in \eqref{tc:eq:M} is linear and contributes nothing.
Taking the removable value of
\eqref{tc:eq:independent-regulators} proves
\eqref{tc:eq:beta-representation}.  Integrable powers at the two
endpoints justify every fixed coefficient extraction.
\end{proof}

All kernels in this formula are explicit finite polynomials.  Indeed,
\begin{equation}\label{tc:eq:logQ}
 \log Q(X,Y)=\sum_{n=2}^{\infty}
  \frac{(-1)^n\zeta(n)}n
      \bigl((X+Y)^n-X^n-Y^n\bigr).
\end{equation}
For instance, with $L=\log(1-t)$ and $M=\log t$,
\begin{align}
 [XY]Qe^{XL+YM}&=LM+\zeta(2),\label{tc:eq:kernel22}\\
 [X^2Y]Qe^{XL+YM}&=\tfrac12L^2M+\zeta(2)L-\zeta(3),\\
 [X^2Y^2]Qe^{XL+YM}
  &=\tfrac14L^2M^2+\zeta(2)LM-\zeta(3)(L+M)
      +\tfrac32\zeta(4)+\tfrac12\zeta(2)^2.
\end{align}
These provide ordinary log-Gamma integral identities without any
endpoint finite-part prescription.

\subsection{The finite composition sum}

Our one-variable multiple polylogarithm convention is
\begin{equation}\label{tc:eq:MPL}
 \operatorname{Li}_{a_0,a_1,\ldots,a_d}(z)
  =\sum_{n_0>n_1>\cdots>n_d\ge1}
      \frac{z^{n_0}}{n_0^{a_0}n_1^{a_1}\cdots n_d^{a_d}},
 \qquad |z|<1.
\end{equation}
Only the outer index receives $z$.  Its branch on
$\mathbb C\setminus[1,\infty)$ is continued from zero by the
usual iterated integrals; there is no branch singularity at zero.
For $r,s\ge1$, $k=r+s$, define the finite combination
\begin{equation}\label{tc:eq:P}
 \boxed{\displaystyle
 P_{r,s}(z)=\sum_{d=1}^{s}\binom{k-d-1}{r-1}
      \sum_{\substack{p_1+\cdots+p_d=k-1\\p_i\ge1}}
                \operatorname{Li}_{1,p_1,\ldots,p_d}(z).}
\end{equation}
Every term has weight $k$ and depth at most $s+1$.

\begin{lemma}[Exact logarithmic face]
\label{tc:lem:logarithm}
If $x>0$ and $B,C\ge0$, set
$q=(C-B)/(x^2+C)$.  Then
\begin{equation}\label{tc:eq:logarithmic-face}
 I_{r+s}\bigl(U\mapsto-\log(x^2+U);B^r,C^s\bigr)
     =(-1)^{r+s-1}r!s!P_{r,s}(q).
\end{equation}
\end{lemma}

\begin{proof}
First suppose $C\ge B$, so $0\le q<1$.  The logarithm is
\[
 -\log(x^2+C)-\log\bigl(1-q(1-t)\bigr).
\]
The constant contributes only a linear term to
$\mathcal M_F(X,Y)$.  Applying the beta moment to the convergent
logarithmic series gives
\begin{equation}\label{tc:eq:log-generating-function}
 \mathcal M_F(X,Y)+(X+Y)\log(x^2+C)
   =X\sum_{m=1}^{\infty}\frac{q^m}{m}
            \prod_{j=1}^{m-1}\frac{1+X/j}{1+(X+Y)/j}.
\end{equation}
This series and its parameter derivatives converge normally for
$|q|<1$ and $X,Y$ in a sufficiently small fixed neighborhood of
zero.  One may use the Gamma-ratio expression for the product to
obtain a polynomial bound in $m$, uniformly on compact parameter
sets; the geometric factor $q^m$ then dominates every derivative.

For a finite product,
\[
 \prod_{j=1}^{m-1}\frac{1+X/j}{1+(X+Y)/j}
   =\prod_{j=1}^{m-1}\left(1-\frac{Y}{j+X+Y}\right).
\]
Choose $d$ nonconstant factors and expand their denominators.
The coefficient of $X^{r-1}Y^s$ is
\begin{align}
 &(-1)^{k-1}\sum_{d=1}^{s}\binom{k-d-1}{r-1}
  \sum_{\substack{p_1+\cdots+p_d=k-1\\p_i\ge1}}
  \sum_{m>j_1>\cdots>j_d\ge1}
                \frac{1}{j_1^{p_1}\cdots j_d^{p_d}}.
                                                     \label{tc:eq:coefficient}
\end{align}
For clarity, the $d$ selected factors supply $(-Y)^d$.
Their denominator expansions supply a total power
$(X+Y)^{k-d-1}$ and sign $(-1)^{k-d-1}$.
The remaining coefficient is precisely
$\binom{k-d-1}{r-1}$.  This proves the sign and every binomial
coefficient in \eqref{tc:eq:coefficient}; no relation among zeta
values has been used.  Summing in $m$ identifies exactly
\eqref{tc:eq:P}.  The coefficient identity from
Theorem~\ref{tc:thm:beta} proves
\eqref{tc:eq:logarithmic-face} for $0\le q<1$.

Both sides are holomorphic in $q$ on
$\mathbb C\setminus[1,\infty)$.  On the integral side this follows
by subtracting the endpoint chord of $-\log(1-q(1-t))$;
on each compact subset the subtracted numerator is uniformly
$O(t(1-t))$.  The identity theorem therefore includes $q<0$,
which is the case $C<B$.  In particular, the analytic continuation
uses real logarithms and real multiple polylogarithms on the whole
real interval $q<1$.
\end{proof}

\subsection{Gamma faces, convergence, and depth}

For $a>0$ and $U\ge0$ define
\begin{equation}\label{tc:eq:G}
 G_a(U)=\log\Gamma(a+i\sqrt U)+\log\Gamma(a-i\sqrt U)
                      -\log(2\pi),
\end{equation}
where the two logarithms are continued from positive real arguments.
Their sum is real.  The classical Gamma product
\cite{dlmfgamma} gives
\begin{equation}\label{tc:eq:Gamma-product}
 G_a(U)=2\log\Gamma(a)-\log(2\pi)
          -\sum_{n=0}^{\infty}
                    \log\left(1+\frac{U}{(n+a)^2}\right).
\end{equation}

\begin{theorem}[Every two-cluster Gamma face]
\label{tc:thm:Gamma}
For $a>0$, $B,C\ge0$ and $r,s\ge1$,
\begin{equation}\label{tc:eq:Gamma-dictionary}
 \boxed{\displaystyle
 I_{r+s}(G_a;B^r,C^s)
   =(-1)^{r+s-1}r!s!
      \sum_{n=0}^{\infty}
       P_{r,s}\!\left(\frac{C-B}{(n+a)^2+C}\right).}
\end{equation}
The spectral sum is ordinary and absolutely convergent, locally
normally in the stated real parameters.  The same identity has a
dictionary of depth at most $\min(r,s)+1$, obtained if necessary by
exchanging the two clusters and using the continued branch on
the negative real axis.  This is an upper bound on a supplied
representation, not a claim of minimal functional depth.
\end{theorem}

\begin{proof}
Apply the ordinary face functional termwise to
\eqref{tc:eq:Gamma-product}.  This is justified directly in the
scalar beta formula: the endpoint-chord remainder of the $n$th
logarithm is bounded in $C^0$ by a constant times
$t(1-t)(C-B)^2(n+a)^{-4}$ on each compact parameter set.
The fixed logarithmic beta kernels are integrable.  Alternatively,
the second derivatives of the log terms converge normally and the
face is continuous in the $C^2$ norm.  Constants disappear, and
Lemma~\ref{tc:lem:logarithm} gives the displayed result.

Each $P_{r,s}(z)$ is $O(z^2)$ at zero, since its first index is one
and its depth is at least two.  The arguments in the tail are
$O(n^{-2})$.  Hence the spectral summands are $O(n^{-4})$,
uniformly on compact subsets.  A finite number of initial arguments
may lie to the left of $-1$ if $C<B$; their analytically continued
values are regular.  This proves absolute and local normal
convergence without a small-displacement hypothesis.

The first expression has depth at most $s+1$.
The same face with the two clusters exchanged has depth at most
$r+1$.  Taking the smaller establishes the asserted bound.
\end{proof}

The symmetry just used gives the precise transformation
\begin{equation}\label{tc:eq:Landen}
 P_{r,s}(z)=P_{s,r}\!\left(\frac{z}{z-1}\right),
       \qquad z\in\mathbb C\setminus[1,\infty).
\end{equation}
It also follows from the classical composition form of the
multiple-polylogarithm Landen connection formula
\cite{okudaueno}.  Here it provides a check on branch conventions
and a useful way to choose the smaller-depth expression.

The first collision coefficient checks all factorial normalizations:
\begin{equation}\label{tc:eq:collision-check}
 I_{r+s}(G_a;B^r,(B+\delta)^s)
  =\frac{(-1)^{r+s-1}}2rs(r+s-2)!\,
     \delta^2\sum_{n=0}^{\infty}
                     \bigl((n+a)^2+B\bigr)^{-2}
       +O(\delta^3).
\end{equation}
This agrees with the predecessor's factorial-variance monomial
formula.  The full dictionary is global for $B,C\ge0$, whereas
that local collision expansion need not converge at the desired
displacement.

\subsection{Two equal displaced masses and a worked ordinary integral}

For $s=1$, \eqref{tc:eq:P} becomes
$P_{r,1}=\operatorname{Li}_{1,r}$ and recovers precisely the
incoming one-displaced theorem.  For $s=2$ it gives the new
all-multiplicity specialization
\begin{equation}\label{tc:eq:two-displaced}
 \boxed{\displaystyle
 P_{r,2}(z)=r\operatorname{Li}_{1,r+1}(z)
     +\sum_{j=1}^{r}\operatorname{Li}_{1,j,r+1-j}(z).}
\end{equation}
Thus $r$ equal base masses and two equal displaced masses have
depth at most three for every $r$.  No increase in the base
multiplicity increases this bound.

For example, set
$E_a(t)=G_a(B+(C-B)t)-(1-t)G_a(B)-tG_a(C)$ and
$q_n=(C-B)/((n+a)^2+C)$.  The case $r=s=2$ gives the completely
ordinary identity
\begin{align}
 &\int_0^1\frac{E_a(t)}{t(1-t)}
                 \bigl(\log t\log(1-t)+\zeta(2)\bigr)\,dt
       \notag\\
 &\qquad=-\sum_{n=0}^{\infty}
   \left(2\operatorname{Li}_{1,3}(q_n)
         +\operatorname{Li}_{1,1,2}(q_n)
         +\operatorname{Li}_{1,2,1}(q_n)\right).
                                                     \label{tc:eq:worked-integral}
\end{align}
At $a=2.3$, $B=0.2$, $C=0.9$, the associated face
$I_4(G_a;B,B,C,C)$ is
\[
 -0.0836254300916287651503714290194933682217772138996833\ldots.
\]
This decimal is an independently checked diagnostic; the identity
is proved above for the full parameter domain.

\subsection{Mass derivatives and anchored polylogarithmic primitives}

Delete the first index one in \eqref{tc:eq:P} and put
\begin{equation}\label{tc:eq:Rpolylog}
 R_{r,s}(z)=\sum_{d=1}^{s}\binom{r+s-d-1}{r-1}
   \sum_{\substack{p_1+\cdots+p_d=r+s-1\\p_i\ge1}}
                       \operatorname{Li}_{p_1,\ldots,p_d}(z).
\end{equation}
The defining strict series gives
$P'_{r,s}(z)=R_{r,s}(z)/(1-z)$.  It follows that
\begin{align}
 \partial_C I_{r+s}(G_a;B^r,C^s)
   &=(-1)^{r+s-1}r!s!
       \sum_{n=0}^{\infty}
           \frac{R_{r,s}(q_n)}{(n+a)^2+C},\label{tc:eq:Cderivative}\\
 \partial_B I_{r+s}(G_a;B^r,C^s)
   &=-(-1)^{r+s-1}r!s!
       \sum_{n=0}^{\infty}
           \frac{R_{r,s}(q_n)}{(n+a)^2+B}.
                                                     \label{tc:eq:Bderivative}
\end{align}
These are total derivatives of the repeated cluster parameters:
$\partial_C$ moves all $s$ copies of $C$ together.  They should not
be confused with a derivative in just one independently labelled mass.
The displayed expressions have depth at most $s$; exchanging clusters
gives a representation with depth at most $\min(r,s)$.

Termwise differentiation is legitimate.  For the tail, the arguments
lie in one fixed small disk, where $P(z)=O(z^2)$ and $R(z)=O(z)$;
local holomorphic extensions and Cauchy estimates give normal
convergence of every fixed parameter derivative.  The finitely many
initial summands stay off their cuts.  Consequently these equations
also give higher polygamma identities after further differentiation.
Since $I_{r+s}(G_a;B^{r+s})=0$, integration of
\eqref{tc:eq:Cderivative} from $C=B$ to the desired endpoint is an
anchored antiderivative with no undetermined constant.

\subsection{Verification and the exact scope of the resolved question}

The supplied script compares two independently organized exact
coefficient calculations: the incoming monomial-face formula at
$0^r,1^s$, and \eqref{tc:eq:P} after the substitution
$z=u/(1+u)$.  It checks $r,s\in\{1,2,3,4\}$ and monomial degrees
$2$ through $15$, together with the corresponding Landen
coefficients, for a total of $448$ exact rational assertions.
Four numerical comparisons at $90$ decimal digits use the ordinary
beta integral on one side and strict harmonic sums with a
Hurwitz-zeta spectral tail on the other.  Their absolute residuals
range from $9.8\times10^{-77}$ to $3.3\times10^{-75}$.
One case has $C-B>a^2+B$, outside the predecessor's local
collision-series disk.  These are floating-point diagnostics,
not certified interval enclosures or proof-assistant proofs.

The result resolves the request at the end of the incoming
spectral-face section for an explicit dictionary when the mass set
has exactly two distinct values, allowing every multiplicity.  It
partially advances Q4 of that report.  It does not evaluate two
\emph{unequal} displaced masses above a third base mass, does not
prove that every two-cluster expression reduces to depth two, and
does not establish a minimal depth theorem.  A precise next target
is the three-cluster case $B^r,C,D$ with $C\ne D$, where an
aggregation into three independent regulator sums should retain
the same ordinary face normalization.

% END INCLUDED SOURCE: sections/04_clusters

% BEGIN INCLUDED SOURCE: sections/05_parity
\section{An exact depth-two example with two Gaussian colors}
\label{par:sec:two-colors}

The Family 2 discussion in the manuscript \cite{repo} concerns
the left-hand side below. We evaluate it in a smaller set of
classical constants by an explicit exact certificate. Our strict convention is
\begin{equation}\label{par:eq:convention}
 \Li_{a,b}(u,v)=\sum_{n>m>0}\frac{u^n v^m}{n^a m^b}.
\end{equation}
In particular, both colors in $\Li_{a,b}(i,i)$ are attached to
summation indices: its numerator is $i^{n+m}$.
This is different from the one-variable value $\Li_{a,b}(i,1)$.
Let $G=\sum_{k\ge0}(-1)^k/(2k+1)^2$ be Catalan's constant.

\begin{theorem}[A weight-three two-color reduction]\label{par:thm:two-colors}
In the convention~\eqref{par:eq:convention},
\begin{equation}\label{par:eq:example}
 \boxed{\displaystyle
 \operatorname{Re}\!\left(\Li_{1,2}(i,i)-\Li_{2,1}(i,i)\right)
 =\frac{9}{32}\zeta(3)-\frac{\pi G}{4}
      +\frac{7\pi^2\log2}{96}.}
\end{equation}
The proof below is a finite rational certificate using double-shuffle
and distribution identities. It does not use numerical recognition or
an unproved inversion formula.
\end{theorem}

\subsection{Convergent relations, including the leading-one limit}

All the series used below converge in their ordinary outer-index order.
Indeed, every $\Li_{2,1}(u,v)$ with $|u|=|v|=1$ converges absolutely,
since $\sum_{n\ge1}H_{n-1}/n^2<\infty$.
For $u\ne1$, $|u|=|v|=1$, write
\[
 \sum_{m<n}\frac{v^m}{m^2}=\Li_2(v)+O(n^{-1}).
\]
Then $\Li_{1,2}(u,v)$ is the sum of
$\Li_2(v)\sum_{n\ge1}u^n/n$ and an absolutely convergent remainder.
The single leading-one values below also have colors different from one.

For $x,y$ among the fourth roots of unity, define
\begin{align}\label{par:eq:R}
 \mathcal R(x,y)={}&
 \Li_{1,2}(x,y/x)+\Li_{2,1}(x,y/x)
       +\Li_{2,1}(y,x/y)\notag\\
 &-\Li_{1,2}(x,y)-\Li_{2,1}(y,x)-\Li_3(xy).
\end{align}
When $x=1$, the two identical leading-one terms in this expression
are canceled \emph{before} evaluation. Thus the definition is then
\begin{equation}\label{par:eq:R-limit}
 \mathcal R(1,y)=
 \Li_{2,1}(1,y)+\Li_{2,1}(y,y^{-1})
       -\Li_{2,1}(y,1)-\Li_3(y).
\end{equation}

\begin{lemma}\label{par:lem:R}
Every value $\mathcal R(x,y)$ used here is zero, including $x=1$.
\end{lemma}

\begin{proof}
For $x\ne1$, the strict stuffle product gives
\begin{equation}\label{par:eq:stuffle}
 \Li_1(x)\Li_2(y)
 =\Li_{1,2}(x,y)+\Li_{2,1}(y,x)+\Li_3(xy).
\end{equation}
One may obtain it by multiplying finite sums with the same cutoff and
partitioning the square into $n>m$, $n<m$, and $n=m$, then using the
convergence just established.
For the shuffle product, use $\omega_0=\dd t/t$ and
$\omega_c=c\,\dd t/(1-ct)$ for $c\ne0$.
Define the endpoint-one iterated integral with $c_1$ outermost by
\[
 J(c_1,\ldots,c_w)=
 \int_{0<t_w<\cdots<t_1<1}\prod_{j=1}^w\omega_{c_j}(t_j).
\]
Expanding the kernels below endpoint one and then taking the convergent
endpoint limit gives
\[
 \Li_{1,2}(u,v)=J(u,0,uv),\qquad
 \Li_{2,1}(u,v)=J(0,u,uv).
\]
The three interleavings of $J(x)$ and $J(0,y)$ therefore give
\begin{equation}\label{par:eq:shuffle}
 \Li_1(x)\Li_2(y)
 =\Li_{1,2}(x,y/x)+\Li_{2,1}(x,y/x)
       +\Li_{2,1}(y,x/y).
\end{equation}
Subtracting~\eqref{par:eq:stuffle} proves the claim for $x\ne1$.

Here is a direct justification of the limiting relation at $x=1$;
no divergent series is assigned a value. Apply the two product formulas
with $x=r\in(0,1)$ and keep $|y|=1$ fixed. The difference of their
leading-one terms satisfies
\begin{align*}
 &\Li_{1,2}(r,y/r)-\Li_{1,2}(r,y)\\
 &\quad=\sum_{m,k\ge1}
       \frac{y^m r^k(1-r^m)}{(m+k)m^2},\\
 &\left|\Li_{1,2}(r,y/r)-\Li_{1,2}(r,y)\right|
 \le -\log(1-r)\bigl(\zeta(2)-\Li_2(r)\bigr)
 =O\!\left((1-r)\log^2(1-r)\right).
\end{align*}
The bound tends to zero. For example,
$\zeta(2)-\Li_2(r)=\int_r^1-\log(1-t)\,\dd t/t
\le (1-r)(1-\log(1-r))/r$.
Each remaining depth-two term is dominated by
$\sum_{n>m}n^{-2}m^{-1}$, and the trilogarithm converges absolutely.
Letting $r\uparrow1$ proves~\eqref{par:eq:R-limit}.
In particular $\mathcal R(1,1)=0$ is Euler's identity
$\Li_{2,1}(1,1)=\zeta(3)$, with its proof included in this limit argument.
\end{proof}

For $u,v\in\{1,-1\}$ choose $\sqrt1=1$ and $\sqrt{-1}=i$, and put
\begin{align}
 \mathcal D_s(u)
 &:=\Li_s(u)-2^{s-1}
       \bigl(\Li_s(\sqrt u)+\Li_s(-\sqrt u)\bigr),
       \label{par:eq:Dsingle}\\
 \mathcal D_{a,b}(u,v)
 &:=\Li_{a,b}(u,v)-2^{a+b-2}
       \sum_{\epsilon,\eta\in\{1,-1\}}
          \Li_{a,b}(\epsilon\sqrt u,\eta\sqrt v).
       \label{par:eq:Ddouble}
\end{align}
Only convergent instances occur in the table below: $s=3$ with $u=\pm1$,
$s=2$ with $u=1$, $s=1$ with $u=-1$, $(a,b)=(2,1)$ with the listed
colors, and $(a,b)=(1,2)$ with $u=-1$.
All these $\mathcal D$ values are zero. To prove this directly, sum the
signs in a finite outer-cutoff sum. They select the even summation indices;
replacing each index by twice a new index gives the factor $2^{d-w}$
at depth $d$ and weight $w$. The outer cutoff becomes its integer half.
Pass to the limit using the preceding convergence checks.

\subsection{The complete rational certificate}

Write $\delta_s=\Li_s(i)-\Li_s(-i)$ for $s=1,2$ and define
\begin{equation}\label{par:eq:target}
 \mathcal T=
 96\bigl(\Li_{1,2}(i,i)-\Li_{2,1}(i,i)\bigr)
 -27\Li_3(1)-24\delta_1\delta_2
 +42\Li_2(1)\Li_1(-1).
\end{equation}
The following table specifies a finite formal identity:
\begin{equation}\label{par:eq:certificate}
 \operatorname{Re}\mathcal T
   =\operatorname{Re}\sum_{j=1}^{15}c_j\mathcal E_j.
\end{equation}
\begin{center}
\begin{tabular}{r r l}
\hline
$j$ & $c_j$ & $\mathcal E_j$ \\
\hline
1 & $-15$ & $\mathcal D_3(1)$ \\
2 & $6$ & $\mathcal D_3(-1)$ \\
3 & $-24$ & $\mathcal R(1,-1)$ \\
4 & $60$ & $\mathcal D_{2,1}(-1,-1)$ \\
5 & $96$ & $\mathcal R(i,1)$ \\
6 & $192$ & $\mathcal R(i,-1)$ \\
7 & $12$ & $\mathcal D_{2,1}(1,1)$ \\
8 & $12$ & $\mathcal R(1,1)$ \\
9 & $-168$ & $\Li_1(i)\mathcal D_2(1)$ \\
10 & $90$ & $\mathcal D_{2,1}(1,-1)$ \\
11 & $96$ & $\mathcal R(i,i)$ \\
12 & $192$ & $\mathcal R(1,i)$ \\
13 & $12$ & $\mathcal D_{1,2}(-1,-1)$ \\
14 & $42$ & $\mathcal D_{1,2}(-1,1)$ \\
15 & $-12$ & $\Li_2(-1)\mathcal D_1(-1)$ \\
\hline
\end{tabular}
\end{center}

For a completely explicit check of~\eqref{par:eq:certificate}, expand
the definitions of $\mathcal R$ and $\mathcal D$, and expand the products
by~\eqref{par:eq:stuffle}. Identify the real part of each coordinate with
the real part of its complex conjugate. With
\[
 X_{a,b}^{p,q}=\operatorname{Re}\Li_{a,b}(i^p,i^q),
 \qquad Z_p=\operatorname{Re}\Li_3(i^p),
\]
both sides of~\eqref{par:eq:certificate} become precisely
\begin{equation}\label{par:eq:expanded-certificate}
 48X_{1,2}^{1,1}+48X_{1,2}^{1,3}+42X_{1,2}^{2,0}
 +42X_{2,1}^{0,2}-144X_{2,1}^{1,1}+48X_{2,1}^{1,3}
 +21Z_0-6Z_2.
\end{equation}
This is an identity in a finite rational vector space before evaluating
any period. The companion script independently constructs the six-term
shuffle residuals directly in strict index-and-color coordinates,
expands every row, and checks equality with exact fractions.

\begin{proof}[Proof of Theorem~\ref{par:thm:two-colors}]
Every $\mathcal E_j$ is zero by Lemma~\ref{par:lem:R} or the finite-cutoff
distribution identities, so~\eqref{par:eq:certificate} gives
$\operatorname{Re}\mathcal T=0$.
The elementary single-value evaluations needed to simplify it are
\[
 \Li_1(i)=-\frac{\log2}{2}+\frac{\pi i}{4},\qquad
 \delta_1=\frac{\pi i}{2},\qquad
 \delta_2=2iG,\qquad
 \Li_2(1)=\frac{\pi^2}{6},\qquad
 \Li_1(-1)=-\log2.
\]
Here the first identity follows from $\Li_1(z)=-\log(1-z)$ with the
principal logarithm, and the formula for $\delta_2$ follows by taking the
odd terms of the absolutely convergent dilogarithm series.
Substitution into~\eqref{par:eq:target} yields
\[
 0=96\operatorname{Re}\!\left(\Li_{1,2}(i,i)-\Li_{2,1}(i,i)\right)
     -27\zeta(3)+24\pi G-7\pi^2\log2.
\]
Division by $96$ proves~\eqref{par:eq:example}.
\end{proof}

The ordinary stuffle product supplies, as a further consequence, the
individual real values
\begin{align}
 \operatorname{Re}\Li_{2,1}(i,i)
   &=\frac{15}{64}\zeta(3)-\frac{\pi^2\log2}{32},
   \label{par:eq:individual-21}\\
 \operatorname{Re}\Li_{1,2}(i,i)
   &=\frac{33}{64}\zeta(3)-\frac{\pi G}{4}
       +\frac{\pi^2\log2}{24}.
   \label{par:eq:individual-12}
\end{align}
Indeed, $\Li_1(i)\Li_2(i)=\Li_{1,2}(i,i)+\Li_{2,1}(i,i)+\Li_3(-1)$,
while $\operatorname{Re}\Li_2(i)=-\pi^2/48$ and
$\Li_3(-1)=-3\zeta(3)/4$ by even--odd splitting.
Together with~\eqref{par:eq:example}, these identities determine the sum
and difference of the two real values.

\subsection{Reconciliation with the printed Family 2 expression}
\label{par:sec:family2-bridge}

The manuscript's displayed Family 2 expression
(\src{double:eq:fam2-example}) uses
\[
 L=\log2,\qquad r=\frac{1+i}{2},\qquad
 \mu_2+i\lambda_2=\Li_2(r),\qquad
 h_{a,b}=\operatorname{Re}\Li_{a,b}(i,1).
\]
To prove that its right side agrees with~\eqref{par:eq:example}, we first
evaluate the required $h_{1,2}$ independently of the printed equality.

\begin{lemma}\label{par:lem:h12}
In the strict convention~\eqref{par:eq:convention},
\begin{equation}\label{par:eq:h12}
 h_{1,2}=-\frac{29}{32}\zeta(3)+\frac{\pi G}{4}
                         +\frac{\pi^2L}{96}.
\end{equation}
\end{lemma}
\begin{proof}
Use the same $\delta_s$ and relation rows $\mathcal E_j$ as in
the certificate~\eqref{par:eq:certificate}, and put
\[
 \mathcal T_h=96\Li_{1,2}(i,1)+87\Li_3(1)
                +24\delta_1\delta_2+6\Li_2(1)\Li_1(-1).
\]
A second exact rational identity in those same coordinates is
\begin{align}\label{par:eq:h12-certificate}
 \operatorname{Re}\mathcal T_h
 =\operatorname{Re}\bigl(&-9\mathcal E_1+18\mathcal E_2
 -12\mathcal E_4-96\mathcal E_5-24\mathcal E_9
 +6\mathcal E_{10}\notag\\
 &-12\mathcal E_{13}+6\mathcal E_{14}+12\mathcal E_{15}\bigr)=0.
\end{align}
As before, this is checked just by expanding the displayed definitions,
stuffle products, and conjugate pairs. Every row on the right vanishes
by the already proved convergent identities; the only shuffle row
needed here has first factor $\Li_1(i)$.
Since $\delta_1\delta_2=-\pi G$ and
$\Li_2(1)\Li_1(-1)=-\pi^2L/6$, the left side is
$96h_{1,2}+87\zeta(3)-24\pi G-\pi^2L$.
Solving proves~\eqref{par:eq:h12}. The companion bridge verifier
reconstructs these nine rows and checks the exact residual separately
from the two-color identity.
\end{proof}

The remaining depth-one and logarithmic quantities are
\begin{equation}\label{par:eq:bridge-values}
 \mu_2=\frac{5\pi^2}{96}-\frac{L^2}{8},\qquad
 \lambda_2=G-\frac{\pi L}{8},\qquad
 h_{1,1}=\frac{L^2}{8}-\frac{\pi^2}{32}.
\end{equation}
For completeness, the classical dilogarithm reflection and Landen
identities \cite{okudaueno} give
\[
 2\mu_2=\zeta(2)-\log r\log\bar r,
 \qquad
 \Li_2(i)+\Li_2(\bar r)=-\tfrac12\log^2(1-i).
\]
Use $\log r=-L/2+i\pi/4$ and $\log(1-i)=L/2-i\pi/4$
to obtain the first two evaluations in~\eqref{par:eq:bridge-values}.
These standard dilogarithm identities follow by differentiating and
fixing the constant at zero, with continuation in the principal domains
containing the displayed arguments. Finally,
$\Li_{1,1}(i,1)=\tfrac12\log^2(1-i)$ follows from the shuffle of
two identical one-letter integrals, and gives the third evaluation.

\begin{corollary}[The printed Family 2 equality]\label{par:cor:family2-printed}
The displayed depth-two identity in the manuscript is valid:
\begin{equation}\label{par:eq:family2-printed}
 \operatorname{Re}\!\left(\Li_{1,2}(i,i)-\Li_{2,1}(i,i)\right)
 =3\zeta(3)-4L\mu_2-\pi\lambda_2-4h_{1,1}L+3h_{1,2}.
\end{equation}
\end{corollary}
\begin{proof}
Substituting~\eqref{par:eq:h12} and~\eqref{par:eq:bridge-values} into the
right side gives
\begin{align*}
 &\left(3-\frac{87}{32}\right)\zeta(3)
 +\left(-1+\frac34\right)\pi G
 +\left(-\frac{20}{96}+\frac{12}{96}
          +\frac{12}{96}+\frac3{96}\right)\pi^2L
 +\left(\frac12-\frac12\right)L^3\\
 &\qquad=\frac9{32}\zeta(3)-\frac{\pi G}{4}
                 +\frac{7\pi^2L}{96}.
\end{align*}
This equals the independently proved left side by
Theorem~\ref{par:thm:two-colors}.
\end{proof}

\paragraph{Independent diagnostic.}
For verification against the strict convention, direct differentiation
and integration of~\eqref{par:eq:convention} give
\[
 \operatorname{Re}\Li_{1,2}(i,i)
 =\operatorname{Re}\int_0^1\frac{i\Li_2(-t)}{1-it}\,\dd t,
 \qquad
 \operatorname{Re}\Li_{2,1}(i,i)
 =-\int_0^1\frac{t\log t\log(1+t)}{1+t^2}\,\dd t.
\]
An optional $80$-digit quadrature in the companion script gives
\[
 \operatorname{Re}\!\left(\Li_{1,2}(i,i)-\Li_{2,1}(i,i)\right)
 =0.11751017576784368621941000460386119840621435030463\ldots.
\]
This is a numerical diagnostic; the finite relation and analytic
arguments above prove the identity. No conclusion about $S_6$ or the
revised $S_8$ conjecture follows from this example.

% END INCLUDED SOURCE: sections/05_parity

% BEGIN INCLUDED SOURCE: sections/06_audit_research
\section{Audit, corrections, and integration status}
\label{sec:audit}

\subsection{What was actually checked}

The source anchor is the repository revision \pin. All eleven archives
in the supplied incoming directory at that revision were unpacked and
indexed. The review concentrated on their theorem and question ledgers
and on the passages used by the present arguments. The manuscript
review concentrated on Gaussian depth reduction, the $S_6$ obstruction,
the revised $S_8$ candidate, and the harmonic/Gamma interfaces.
This targeted review is not a certification of all manuscript volumes
or every theorem in the incoming reports.

Three source questions have precise new outcomes:

\begin{enumerate}
\item \textbf{Ordered boundary report, research question 2.}
The positive constants in \src{sections/08_gamma_basis.tex} are
evaluated by Theorem~\ref{gn:thm:K}. Its proof uses the source's
established top-row identities and a classical integral with an
explicitly checked normalization. The matrix $T_m$ and all column
orders are stated before use. The later-index determinant in
\eqref{gn:eq:F-explicit} is a different object.

\item \textbf{Balanced Stieltjes report, research question 4.}
The homogeneous colored balancing, its natural Lerch coordinates,
its ordinary Mellin moments, and its finite root-of-unity filter
are supplied in Sections~\ref{sec:colored}--\ref{sec:color-limit}.
Uniqueness is asserted for coefficients independent of the shift
and for a decay condition across an open interval of shifts.
This is the same kind of hypothesis needed in the predecessor's
uncolored uniqueness theorem.

\item \textbf{Lattice-contact report, end of the spectral-face section
and research question Q4.}
Theorem~\ref{tc:thm:Gamma} treats exactly two distinct mass values at
every pair of positive multiplicities. It supplies a finite
dictionary per spectral term. A problem involving two unequal
displaced masses above a third base value has three clusters and
remains outside this theorem.
\end{enumerate}

The companion source manifest records the archive names and their
SHA-256 hashes. The integration notes list exact predecessor paths
and theorem labels, so a later integration can reconcile the new
results with a changed repository head.

\subsection{Corrections and qualifications}

No false mathematical theorem was identified in the particular
incoming passages used for these three developments. It would be
misleading to manufacture an erratum from an open question that
the source already labels correctly. The following distinctions
are nevertheless essential when integrating the continuation.

\paragraph{The colored generator changes.}
For a nontrivial color, $\Phi(\omega,s,a)$ has no pole at $s=1$.
The appropriate balanced kernel is \eqref{col:kernel}. Copying
the extra uncolored term $\log^{q+1}x/(q+1)$ would spoil its
large-$x$ cancellation. This is a correction to that potential
extension, not a claim that the incoming uncolored definition
is erroneous.

\paragraph{Fixed-shift uniqueness can fail.}
The open-interval clause in Theorem~\ref{col:unique} cannot simply
be deleted on the full color domain. The exceptional-color
example following that theorem exhibits simultaneous vanishing
of consecutive tail coordinates. At such a point, a logarithm
times a lower balanced kernel can preserve the required decay.

\paragraph{Constant terms depend on the chosen expression.}
The raw $\lambda$-constant term of
$M_{m,q}(\lambda,a)$ differs from that of
$e^{-a\lambda}M_{m,q}(\lambda,a)$ by $L_{m,q}a^m/m$.
Equations \eqref{col:renorm} and \eqref{col:raw-CT} retain that
term explicitly. This avoids an incorrect identification of a
color limit with an uncolored ordinary moment.

\paragraph{Positivity does not settle every later determinant.}
The positive density \eqref{gn:eq:density} evaluates the
normalizing determinant. The source's later rows involve
$\Im(x+i/2)^{2r}$ with oscillatory signs. One cannot infer all
later signs by applying positivity of a Vandermonde product to
a different matrix. In particular, no all-pairs nonvanishing
claim follows from the closed formula for $K_m$ alone.

\paragraph{Finite normal moments do not give a finite-dimensional
normal distribution.}
The statistic in Theorem~\ref{gn:thm:gaussian} is bounded at fixed
$n,\sigma$. We prove equality of moments only through degree
$n+1$, and we do not assert that this range is maximal.
The distributional limit in Corollary~\ref{gn:cor:CLT} has a
separate moment-method proof.

\paragraph{Branch and color conventions remain explicit.}
The two-cluster $\Li_{\boldsymbol a}(z)$ assigns its argument only
to the largest summation index. Colored depth-two notation uses
two independent color factors and is defined separately.
When $C<B$, the cluster formula uses the continuation on
$\C\setminus[1,\infty)$; substituting a divergent power series
at a negative argument of modulus greater than one is not valid.
Likewise, the ordinary $\omega$ series for the colored balanced
kernel is used only inside its disk of convergence before
continuation to the full slit domain.

\paragraph{A small editorial correction.}
At the pinned revision, the manuscript README repeats consecutively
the paragraph describing the complete $1{,}032$-input first-letter
$X$ family, its rank $1{,}008$, and its kernel dimension $24$.
The optional patch in the integration directory removes one
copy and leaves the mathematical statements unchanged.

\subsection{Evidence and its limits}

Analytic proofs establish the identities. The accompanying exact
and numerical calculations provide independent checks on the
algebra, normalization, and selected parameter values.

The Gamma normalization suite evaluates the original coefficient
matrices at $m=1,2,3$, the parameterized determinants at unequal
real and complex shifts, and scalar digamma/polygamma insertions.
It contains fourteen checks at $42$ decimal digits; the largest
recorded relative residual is approximately $7.1\times10^{-43}$.
The finite-Gaussian suite constructs exponential jets of ordinary
one-dimensional moment integrals and takes their determinant by
finite polynomial arithmetic. It does not use the paired Gamma
product on the left side. Its fifteen checks include the fourth
moment at $n=3$ and non-quarter shifts; the largest recorded
absolute residual is approximately $4.9\times10^{-41}$.

The cluster suite contains $448$ exact rational assertions,
comparing two different coefficient constructions and checking
cluster exchange. Four independent numerical comparisons at
$90$ digits include a displacement outside the predecessor's
local collision-series disk. The colored suite supplies $41$ exact assertions and eleven
numerical checks at $50$ digits, covering finite tail and zeta
coefficients, a simultaneous-tail-zero example, ordinary moments,
root-of-unity specializations, and an anchored primitive. Its
largest numerical residual is approximately $5.4\times10^{-37}$.
Three further $55$-digit checks evaluate a negative cluster argument
$q=-4$, its Landen transform, and the oriented anchored primitive.
The separate depth-two certificate uses exact rational
relations, with its conventions and limiting argument explained
in the corresponding section. A second certificate, using nine
of the same relation rows, evaluates the intermediate $h_{1,2}$
and reconciles the actual printed Family 2 right side exactly.

The JSON records specify the checks actually performed. Floating
point residuals are diagnostics, not certified error intervals.
An exact certificate verifies an algebraic reduction under its
stated analytic identities; it is not by itself a formalization
of convergence in a proof assistant. No decimal relation search
is promoted to a proof of $S_6$ or $S_8$.

\section{Further research questions}
\label{sec:research}

The results suggest several concrete next problems. The order below
prioritizes exact identities and structural classification over
estimates. Where a question is already in the supplied sources,
the proposed next step is stated relative to the new formulas.

\begin{enumerate}[label=\textbf{Q\arabic*.},leftmargin=*]

\item \textbf{Find the first non-Gaussian moment and its exact defect.}
For the statistic $X_{n,\sigma}$, determine
\[
 \mathbb E X_{n,\sigma}^{\,n+2}
       -\mathbb E Z^{\,n+2},\qquad Z\sim N(0,1),
\]
and decide whether $n+1$ is the exact universal matching degree.
The distinct-parameter proof stops because there are only
$n+1$ pair labels. Repeated powers of one label bring in
$S_3,S_5,\ldots$; the complete generating identity
\eqref{gn:eq:polygamma-germ} gives coupled equations for these
new coordinates. This is an algebraic obstruction to isolate,
not evidence that the next moment must be nonzero in every case.

\item \textbf{Classify joint polygamma moments.}
Give a finite combinatorial rule for the mixed cumulants of
$S_1,S_3,S_5,\ldots$ obtainable from the paired deformation.
The logarithm of the Gamma product involves interactions of at
most two parameter labels, but repeated powers of a label carry
higher polygamma statistics. A useful theorem would state
precisely which joint moments can be separated at dimension $n$,
and give their coefficients as polygamma values at $2\sigma$.
At $\sigma=1/4$, reflection should yield especially short
rational combinations of powers of $\pi$.

\item \textbf{Extend the determinant evaluation to later source rows.}
For the determinant in \eqref{gn:eq:F-explicit}, express a
prescribed family of later row sets as the evaluated Gustafson
normalizer times an explicit symmetric polynomial or a finite
hypergeometric sum. Such an identity would preserve the
oscillatory polynomial insertions exactly. It could address
the source's fixed-gap sign changes without an unjustified
total-positivity argument. The gap-six statement
$\det F_{b,3}>0$ for $b\ge7$ remains a concrete open target
\cite{repoordered}.

\item \textbf{Classify simultaneous zeros of colored tails.}
For
$A_k(\omega,a)=(\omega\partial_\omega+a)^k(1-\omega)^{-1}$,
classify all common zeros of consecutive $A_k$ at $a>0$ and
$\omega\in\Omega$. The example in this article proves that
consecutive zeros do occur. Determine the resulting dimension
of the space of homogeneous combinations with enhanced decay,
and the exact extra ordinary moments that each zero pattern
permits. A classification on the unit circle may be simpler
than one on the entire slit domain.

\item \textbf{Give a character-theoretic exceptional-moment dictionary.}
Apply finite combinations of root-of-unity colors to isolate
Dirichlet characters. At a negative integer where the relevant
$L$-function vanishes, determine the corresponding ordinary
Stieltjes-difference integral and simplify it using the
functional equation. The alternating beta family
\eqref{col:beta-formula} is the first case. A complete statement
should distinguish primitive and imprimitive characters and
retain their conductor and parity factors.

\item \textbf{Evaluate the first bilinear colored moment.}
For positive shifts, determine a natural ordinary domain and an
exact function class for
\[
 \int_0^\infty x^{u-1}
                  C_p^\omega(x,a)C_q^\eta(x,b)\dd x.
\]
The scalar Mellin formula suggests a Mellin convolution, but
the transform of a product is not the product of transforms.
One target is a two-parameter beta/Feynman integral whose
integer contacts reduce to a finite family of multiple-Lerch
coordinates. Coincident shifts and inverse colors are plausible
specializations to analyze first.

\item \textbf{Resolve the three-cluster spectral face.}
For $B^r,C,D$ with three distinct nonnegative masses, aggregate
the independent regulators into three sums before extracting
the squarefree coefficient. The two-cluster proof then suggests
a two-simplex hypergeometric integral and a finite iterated
integral dictionary. Determine the correct multivariable
polylogarithm arguments and branch domain. The answer should
explain what changes when $C=D$, where the depth-three formula
\eqref{tc:eq:two-displaced} must be recovered.

\item \textbf{Determine genuine depth reductions and lattice closures.}
The bound $\min(r,s)+1$ is a supplied representation bound.
Can the two-cluster combinations be reduced further as
functions? Independently, characterize parameters for which
the outer sum in \eqref{tc:eq:Gamma-dictionary} collapses to
finitely many multiple Gamma, Hurwitz, or polylogarithm
coordinates. These are two different problems: reducing
depth termwise does not eliminate an infinite spectral sum.

\item \textbf{Iterate anchored primitives with their constants intact.}
The colored primitive \eqref{col:primitive} lowers the Lerch
order by one; the cluster derivative deletes an initial index
one in the polylogarithm word. Iterate both operations and
identify a canonical Barnes-function normalization for the
root-of-unity cases. At every stage, retain explicit anchors.
Determine whether the normalized color-one limit commutes
with these primitives and identify any additional polynomial
terms when it does not.

\item \textbf{Supply a new relation for $S_6$ and revised $S_8$.}
The manuscript already distinguishes the conjectural period
equalities from its exact bounded-span and shuffle-obstruction
theorems. A proof must introduce a relation outside the
already excluded span, or a new analytic integral identity,
with all branch and endpoint terms included. The weight-three
worked reduction in this article is a useful small exact
certificate; it does not remove the higher-weight obstruction.
An appropriate deliverable is a rational certificate whose
primitive rows have independent analytic proofs, not only an
enlarged numerical relation search.

\item \textbf{Formalize the reusable analytic interfaces.}
The most useful first targets are continuous Cauchy--Binet with
dominated parameter differentiation, finite coefficient
extraction from normally convergent Mellin germs, and the
independent-regulator face identity
\eqref{tc:eq:independent-regulators}. Formalizing these
interfaces would allow many subsequent identities to be
checked by finite algebra while retaining the convergence
arguments that justify that algebra.

\end{enumerate}


% END INCLUDED SOURCE: sections/06_audit_research

% BEGIN INCLUDED SOURCE: references
\clearpage
\begin{thebibliography}{99}
\small\raggedright
\addcontentsline{toc}{section}{References}

\bibitem{repo}
V.~Reshetnikov, \emph{ProveIt: Analysis/Polylogarithms manuscript},
repository revision \pin, especially the Gaussian depth-reduction
and harmonic/Gamma chapters.
\url{https://github.com/VladimirReshetnikov/ProveIt/tree/c4036d60715916ae72dd18983782fdb182840b86/Analysis/Polylogarithms/docs/manuscript}.

\bibitem{repoordered}
ProveIt incoming report,
\emph{Ordered Boundary Identities and Gamma Kernels},
11 October 2026,
archive \src{ProveIt_Ordered_Boundary_Identities_2026-10-11.zip},
inner \src{sections/08_gamma_basis.tex}
(\src{gb:basis}, \src{gb:moment-identities}, \src{gb:normalization})
and \src{sections/09_audit_research.tex}, question 2,
at revision \pin.
\url{https://github.com/VladimirReshetnikov/ProveIt/tree/c4036d60715916ae72dd18983782fdb182840b86/docs/incoming}.

\bibitem{repobalanced}
ProveIt incoming report,
\emph{Balanced Stieltjes Moments: Bernoulli Closure, Harmonic Filters,
and Polylogarithmic Resonances},
11 October 2026,
archive \src{ProveIt_Balanced_Stieltjes_Moments_2026-10-11.zip},
inner \src{article/Balanced_Stieltjes_Moments.tex},
especially \src{bsm:thm:unique}, \src{bsm:thm:master},
\src{bsm:thm:integer}, and research question 4,
at revision \pin.
Available in the incoming directory cited above.

\bibitem{repolattice}
ProveIt incoming report,
\emph{Lattice Contacts and Mixed Zeta Identities},
11 October 2026,
archive
\src{ProveIt_Lattice_Contacts_and_Mixed_Zeta_Identities_2026-10-11.zip},
inner \src{sections/07_spectral_faces.tex}
(\src{sfj:thm:monomial}, \src{sfj:thm:gamma-polylog})
and \src{sections/08_audit_and_research.tex}, Q4,
at revision \pin.
Available in the incoming directory cited above.

\bibitem{gustafson1994}
R.~A. Gustafson,
\emph{Some $q$-beta and Mellin--Barnes integrals on compact Lie groups
and Lie algebras},
Transactions of the American Mathematical Society \textbf{341}
(1994), 69--119.
\url{https://doi.org/10.1090/S0002-9947-1994-1139492-3}.
The normalization used here was checked in the primary reproduction
\cite[equation (2)]{derkachov2017}; no uninspected original theorem
number is asserted.

\bibitem{derkachov2017}
S.~\'E. Derkachov and A.~N. Manashov,
\emph{Spin Chains and Gustafson's Integrals},
Journal of Physics A: Mathematical and Theoretical \textbf{50}
(2017), 294006; arXiv:1611.09593v2, equation (2).
\url{https://doi.org/10.1088/1751-8121/aa749a}.
\url{https://arxiv.org/html/1611.09593v2}.

\bibitem{okudaueno}
J.-I. Okuda and K. Ueno,
\emph{The Sum Formula of Multiple Zeta Values and Connection Problem
of the Formal Knizhnik--Zamolodchikov Equation},
arXiv:math/0310259v2 (2004), especially equation (2).
\url{https://arxiv.org/abs/math/0310259}.

\bibitem{dlmfgamma}
NIST Digital Library of Mathematical Functions,
\S5.8, \emph{Infinite Products}, especially 5.8.3.
\url{https://dlmf.nist.gov/5.8}.

\bibitem{dlmfbeta}
NIST Digital Library of Mathematical Functions,
\S5.12, \emph{Beta Function}, especially Euler's beta integral.
\url{https://dlmf.nist.gov/5.12}.

\bibitem{dlmfpolygamma}
NIST Digital Library of Mathematical Functions,
\S5.15, \emph{Polygamma Functions}.
\url{https://dlmf.nist.gov/5.15}.

\bibitem{dlmfhurwitz}
NIST Digital Library of Mathematical Functions,
\S25.11, \emph{Hurwitz Zeta Function}.
\url{https://dlmf.nist.gov/25.11}.

\bibitem{dlmflerch}
NIST Digital Library of Mathematical Functions,
\S25.14, \emph{Lerch's Transcendent},
especially 25.14.5 and 25.14.3\_1.
\url{https://dlmf.nist.gov/25.14}.

\bibitem{dlmfdirichlet}
NIST Digital Library of Mathematical Functions,
\S25.15, \emph{Dirichlet $L$-Functions}, functional equation
for primitive characters and trivial zeros.
\url{https://dlmf.nist.gov/25.15}.

\end{thebibliography}

% END INCLUDED SOURCE: references

\end{document}
